\documentclass[11pt]{article}
\usepackage[margin=1in]{geometry}
\usepackage{setspace}
\usepackage{amsmath}
\usepackage{amssymb}
\usepackage{amsthm}
\usepackage{booktabs}
\usepackage{graphicx}
\usepackage{longtable}
\usepackage{array}
\usepackage{multirow}
\usepackage{enumitem}
\providecommand{\IfDocumentMetadataT}[1]{}
\providecommand{\IfDocumentMetadataTF}[2]{#2}
\usepackage{hyperref}
\usepackage{natbib}
\usepackage{url}
\usepackage{microtype}
\setlength{\parskip}{0.35em}
\setlength{\parindent}{0pt}
\onehalfspacing
\hypersetup{colorlinks=true,linkcolor=blue,citecolor=blue,urlcolor=blue}
\newcommand{\EFW}{\mathrm{EFW}}
\newcommand{\GDP}{\mathrm{GDP}}
\newcommand{\ITT}{\mathrm{ITT}}
\newcommand{\LATE}{\mathrm{LATE}}
\newcommand{\RD}{\mathrm{RD}}
\newcommand{\IV}{\mathrm{IV}}
\newcommand{\TwoSLS}{\mathrm{2SLS}}
\newcommand{\cov}{\mathrm{Cov}}
\newcommand{\var}{\mathrm{Var}}
\newcommand{\E}{\mathbb{E}}
\newcommand{\1}{\mathbf{1}}
\newcommand{\eps}{\varepsilon}
\newcommand{\mc}{\mathcal}
\newcommand{\clust}{\mathrm{clust}}
\newcommand{\se}{\mathrm{se}}
\newcommand{\sign}{\mathrm{sign}}
\newcommand{\dd}{,\mathrm{d}}
\newcommand{\argmin}{\mathrm{arg,min}}
\newcommand{\argmax}{\mathrm{arg,max}}
\newcommand{\reform}{\mathrm{reform}}
\newcommand{\reversal}{\mathrm{reversal}}
\newcommand{\state}{\mathrm{state}}
\newcommand{\tail}{\mathrm{tail}}
\newcommand{\Hc}{\mathrm{H}}
\newcommand{\Lc}{\mathrm{L}}
\newcommand{\bW}{\mathbf{W}}
\newcommand{\bX}{\mathbf{X}}
\newcommand{\bZ}{\mathbf{Z}}
\newcommand{\bS}{\mathbf{S}}
\newcommand{\btheta}{\boldsymbol{\theta}}
\newcommand{\bbeta}{\boldsymbol{\beta}}
\newcommand{\bgamma}{\boldsymbol{\gamma}}
\newcommand{\btau}{\boldsymbol{\tau}}
\newcommand{\ubar}{\bar{u}}
\newcommand{\cbar}{\bar{c}}
\newcommand{\ybar}{\bar{y}}
\newcommand{\mbar}{\bar{m}}
\newcommand{\RDcut}{0}
\newcommand{\bw}{h}
\newcommand{\bwLR}{r}
\newcommand{\CCT}{\text{CCT}}
\newcommand{\CFT}{\text{CFT}}
\newcommand{\PWT}{\text{PWT}}
\newcommand{\WDI}{\text{WDI}}
\newcommand{\DPI}{\text{DPI}}
\newcommand{\VDEM}{\text{V-Dem}}
\newcommand{\VPARTY}{\text{V-Party}}
\newcommand{\CLEA}{\text{CLEA}}
\newcommand{\ParlGov}{\text{ParlGov}}
\newcommand{\MARPOR}{\text{MARPOR}}
\newcommand{\LV}{\text{Laeven--Valencia}}
\newcommand{\N}{\mathbb{N}}
\newcommand{\R}{\mathbb{R}}
\newcommand{\C}{\mathbb{C}}
\newcommand{\T}{\mathbb{T}}
\newcommand{\Indic}{\mathrm{Indic}}
\newcommand{\Prob}{\mathbb{P}}
\newcommand{\Min}{\mathrm{Min}}
\newcommand{\Max}{\mathrm{Max}}
\newcommand{\logit}{\mathrm{logit}}
\newcommand{\probit}{\mathrm{probit}}
\newcommand{\Bern}{\mathrm{Bernoulli}}
\newcommand{\Normal}{\mathrm{Normal}}
\newcommand{\Unif}{\mathrm{Unif}}
\newcommand{\iid}{\mathrm{iid}}
\newcommand{\as}{\mathrm{a.s.}}
\newcommand{\plim}{\mathrm{plim}}
\newcommand{\To}{\Rightarrow}
\newcommand{\toP}{\xrightarrow{p}}
\newcommand{\toD}{\xrightarrow{d}}
\newcommand{\toas}{\xrightarrow{a.s.}}
\newcommand{\ddx}{\frac{\mathrm{d}}{\mathrm{d}x}}
\newcommand{\ddt}{\frac{\mathrm{d}}{\mathrm{d}t}}
\newcommand{\ddm}{\frac{\mathrm{d}}{\mathrm{d}m}}
\newcommand{\ddh}{\frac{\mathrm{d}}{\mathrm{d}h}}
\newcommand{\given}{,\big|,}
\newcommand{\abs}[1]{\left|#1\right|}
\newcommand{\norm}[1]{\left|#1\right|}
\newcommand{\paren}[1]{\left(#1\right)}
\renewcommand{\brack}[1]{\left[#1\right]}
\newcommand{\curly}[1]{\left{#1\right}}
\newcommand{\anglebr}[1]{\left\langle#1\right\rangle}
\newcommand{\ceil}[1]{\left\lceil#1\right\rceil}
\newcommand{\floor}[1]{\left\lfloor#1\right\rfloor}
\newcommand{\ud}{,\mathrm{d}}
\newcommand{\ubarh}{\underline{h}}
\renewcommand{\hbar}{\overline{h}}
\newcommand{\ubarv}{\underline{v}}
\newcommand{\vbar}{\overline{v}}
\newcommand{\ubarbeta}{\underline{\beta}}
\newcommand{\barbeta}{\overline{\beta}}
\newcommand{\ubarDelta}{\underline{\Delta}}
\newcommand{\barDelta}{\overline{\Delta}}
\newcommand{\ubarpsi}{\underline{\psi}}
\newcommand{\barpsi}{\overline{\psi}}
\newcommand{\ubarphi}{\underline{\phi}}
\newcommand{\barphi}{\overline{\phi}}
\newcommand{\ubarpi}{\underline{\pi}}
\newcommand{\barpi}{\overline{\pi}}
\newcommand{\ubaralpha}{\underline{\alpha}}
\newcommand{\baralpha}{\overline{\alpha}}
\newcommand{\ubarY}{\underline{Y}}
\newcommand{\barY}{\overline{Y}}
\newcommand{\ubarE}{\underline{E}}
\newcommand{\barE}{\overline{E}}
\newcommand{\ubarS}{\underline{S}}
\newcommand{\barS}{\overline{S}}
\newcommand{\ubarD}{\underline{D}}
\newcommand{\barD}{\overline{D}}
\newcommand{\ubarZ}{\underline{Z}}
\newcommand{\barZ}{\overline{Z}}
\newcommand{\ubarX}{\underline{X}}
\newcommand{\barX}{\overline{X}}

\begin{document}

%%%%%%%%%%%%%%%%%%%%%%%%%%%%%%%%%%%%%%%%%%%%%%%%%%%%%%%%%%%%%%%%%%%%%%%%%%%%%%%
% Targeted web-confirmed factual anchors used in this proposal:
% EFW annual report: coverage, components, download/free access, version notes. ([Fraser Institute][1])
% CLEA coverage and versioning (Oct 15, 2025 update) and public-good access notes. ([electiondataarchive.org][2])
% ParlGov scope (EU/OECD, 1900--2023) and access points. ([parlgov.org][3])
% DPI scope (1975--2020, 180 countries, includes ideology measures) ([Inter-American Development Bank][4])
% V-Party dataset access and scope summary (QoG + V-Dem pages). ([v-dem.net][5])
% Manifesto Project API/key requirement (manifestoR doc). ([CRAN][6])
% WDI API access and licensing (CC BY 4.0) ([datahelpdesk.worldbank.org][7])
% PWT licensing/coverage (10.0 page) and new release (11.0 through 2023) ([University of Groningen][8])
% Systemic Banking Crises database update coverage (1970--2011). ([SSRN][9])
% RD inference foundations (CCT 2014; CFT 2015; Lee 2008). ([rdpackages.github.io][10])
%%%%%%%%%%%%%%%%%%%%%%%%%%%%%%%%%%%%%%%%%%%%%%%%%%%%%%%%%%%%%%%%%%%%%%%%%%%%%%%

\begin{center}
{\LARGE \textbf{Close Elections, Economic Freedom, and Macro Performance:}}\
{\LARGE \textbf{A Cross-Country RD Spine and a Nonlinear RD-IV Agenda}}\
\vspace{0.3em}
{\large Publication-grade proposal (identification-first, referee-proof)}\
\vspace{0.3em}
{\large January 25, 2026}\
\end{center}

\vspace{0.8em}

%%%%%%%%%%%%%%%%%%%%%%%%%%%%%%%%%%%%%%%%%%%%%%%%%%%%%%%%%%%%%%%%%%%%%%%%%%%%%%%
\section*{0) Executive Pitch}
\textit{(Constraint: $\leq$ 1 page; concise pitch + most defensible takeaway.)}

\vspace{0.4em}

\textbf{Title (tight + specific).}\
Close Elections as Quasi-Random Market Orientation:\
Estimating the Local Impact of Election-Induced Economic Freedom Changes on Growth and Tail Risk.

\vspace{0.4em}

\textbf{Abstract (150--200 words).}\
This project builds a cross-country regression discontinuity (RD) design using close elections as
quasi-random assignment of a more-market versus less-market competitor.
Part 1 establishes an `improved RD'' spine: robust bias-corrected local polynomial RD
(and complementary local randomization inference) to estimate the local causal effect of electing a
more-market winner on subsequent changes in the Economic Freedom of the World (EFW) index and its
components, and on macro outcomes (growth, inflation, investment, tail risk and crisis/drawdown
metrics) across horizons aligned to electoral terms.
Part 2 leverages the Part-1 first stage to identify the effect of EFW on macro outcomes via
RD-IV (a fuzzy RD / local Wald strategy), with explicit honesty about exclusion failures:
winners plausibly affect macro outcomes through channels not fully captured by EFW.
Accordingly, we treat `EFW $\rightarrow$ macro'' as a local causal claim under additional
assumptions, and we complement point estimates with sensitivity analysis and bounds for plausible
direct winner effects.
A central innovation is a pre-specified framework to detect and interpret nonlinearities, asymmetries
(reforms vs reversals), and state dependence (baseline institutions, inflation regime, and initial EFW)
in the macro response to EFW changes induced by close elections.

\vspace{0.4em}

\textbf{Why it matters (5 bullets).}
\begin{itemize}[leftmargin=1.5em]
\item Many macro-policy debates hinge on whether `market institutions'' causally shape growth and
macro stability, yet cross-country evidence is often confounded by selection and reverse causality.
\item EFW is widely used as a summary measure of market-supporting institutions, but the causal
meaning of EFW changes remains contested and under-identified in macro settings. :contentReference[oaicite:10]{index=10}
\item Close elections offer a rare quasi-experimental handle on institutional change at national scale.
\item Tail outcomes (crises, drawdowns, inflation spikes) are first-order welfare events not captured
by average growth regressions.
\item The policy relevance is highest when effects are nonlinear and state-dependent (`when do reforms
matter, and when do reversals hurt most?'').

\end{itemize}

\textbf{What is new (5 bullets).}
\begin{itemize}[leftmargin=1.5em]
\item A cross-country `improved RD'' election design focused on identifying a \emph{first stage} from
close-election winners to \emph{subsequent EFW movements} (overall and components), not merely winner
effects on spending or transfers.
\item Term-aligned horizon mapping that separates immediate policy shifts from longer-run dynamics,
with explicit handling of overlapping elections as a contamination problem.
\item A Part-2 RD-IV agenda that directly targets \emph{EFW $\rightarrow$ macro} and does not hide the
exclusion problem; it adds sensitivity/bounds as first-class outputs rather than afterthoughts.
\item A pre-specified nonlinear/asymmetric/state-dependent estimand set for EFW effects (levels,
changes, tails), avoiding post-hoc `winner ideology'' narratives.
\item Component-wise analysis that confronts the fact that some EFW subcomponents overlap mechanically
with certain macro outcomes (e.g., inflation inside ``Sound Money''), and redesigns outcomes/indices
accordingly. ([Fraser Institute][1])

\end{itemize}

\textbf{Why identification is credible (5 bullets, with explicit caveats).}
\begin{itemize}[leftmargin=1.5em]
\item \textbf{Part 1 credibility:} Near the zero-margin cutoff, winner assignment is plausibly as-good-as
random conditional on smooth potential outcomes (standard close-election RD logic). ([ScienceDirect][11])
\item \textbf{Modern RD inference:} Default estimation uses robust bias-corrected local polynomial RD and
MSE-optimal bandwidths with valid inference; plus local randomization robustness. ([rdpackages.github.io][10])
\item \textbf{Diagnostics:} density/manipulation tests, covariate balance, donut RD, placebo cutoffs and
pre-trend outcomes, and system-specific falsifications.
\item \textbf{Caveat 1 (orientation):} More-market vs less-market labeling is imperfect; we pre-specify
a ranked ideology-scoring pipeline and report robustness across label sources.
\item \textbf{Caveat 2 (Part 2 exclusion):} Election winners may affect macro outcomes through channels
not captured by EFW; thus, ``EFW $\rightarrow$ macro'' requires additional assumptions and will be
accompanied by bounds and sensitivity analysis, not presented as unconditional truth.

\end{itemize}

\textbf{Most defensible takeaway (one paragraph).}\
The paper’s most defensible claim is reduced-form:
\emph{electing a more-market winner in a close election causes measurable changes in EFW (first stage)
and causes shifts in macro outcomes and tail risks over the subsequent term, locally at the margin of
electoral victory.}
A stronger ``EFW $\rightarrow$ macro'' interpretation is offered only as a local causal claim under
explicitly stated exclusion-type assumptions and is backed by sensitivity/bounding exercises that
quantify how large non-EFW winner channels would need to be to overturn conclusions.

\vspace{1.0em}

%%%%%%%%%%%%%%%%%%%%%%%%%%%%%%%%%%%%%%%%%%%%%%%%%%%%%%%%%%%%%%%%%%%%%%%%%%%%%%%
\section*{1) Research Questions and Estimands}
\textit{(Explicit, conservative, and estimand-first.)}

\vspace{0.4em}

\textbf{Notation.}\
Countries indexed by $c \in \mc{C}$.\
Elections indexed by $e \in \mc{E}_c$.\
Let $t_{ce}$ be the calendar year of election $e$ in country $c$.\
Let $m_{ce}$ be the (signed) margin of victory for the more-market competitor (defined below).\
Treatment indicator $D_{ce} = \1{m_{ce} > 0}$ (more-market competitor wins).\
Cutoff at $m=0$.

\vspace{0.3em}

Let $\EFW_{c,t}$ be the EFW index for country $c$ in year $t$ (overall or a component).\
EFW is available for up to 165 jurisdictions back to 1970, with 45 components in 5 areas. ([Fraser Institute][1])\
(Important measurement note: EFW components include ``Sound Money'' subcomponents that incorporate
inflation; we treat such overlap carefully in outcome definitions.)

\vspace{0.4em}

Let $Y_{c,t}$ denote a macro outcome (e.g., real GDP per capita growth, inflation, investment share).\
Let $Y_{ce}(h)$ denote an election-horizon outcome defined below.

\vspace{0.8em}

%%%%%%%%%%%%%%%%%%%%%%%%%%%%%%%%%%%%%%%%%%%%%%%%%%%%%%%%%%%%%%%%%%%%%%%%%%%%%%%
\subsection*{1.1 Part 1 estimands}
\textit{(RD reduced form; what we can claim under minimal RD assumptions.)}

\vspace{0.4em}

\textbf{Primary Part-1 estimand family: treatment-on-outcome RD discontinuities.}\
For any election-level outcome $O_{ce}$ (constructed from post-election time series),
define the RD estimand at the cutoff:
\begin{equation}
\tau_O
=
\lim_{m \downarrow 0} \E\big[ O_{ce} \given m_{ce}=m \big]
-
\lim_{m \uparrow 0} \E\big[ O_{ce} \given m_{ce}=m \big].
\end{equation}

\vspace{0.4em}

\textbf{EFW first-stage estimands (overall + components).}\
Define horizon-$h$ change in EFW relative to a pre-election baseline year:
\begin{equation}
\Delta \EFW_{ce}(h)
=
\EFW_{c,t_{ce}+h}
-
\EFW_{c,t_{ce}-1}.
\end{equation}

\textbf{Primary EFW estimands:}
\begin{equation}
\tau_{\Delta \EFW(h)}
=
\lim_{m \downarrow 0} \E\big[\Delta \EFW_{ce}(h)\given m_{ce}=m\big]
-
\lim_{m \uparrow 0} \E\big[\Delta \EFW_{ce}(h)\given m_{ce}=m\big].
\end{equation}

We estimate $\tau_{\Delta \EFW(h)}$ for:
\begin{itemize}[leftmargin=1.5em]
\item overall EFW index,
\item the 5 EFW areas (Size of Government; Legal System \& Property Rights; Sound Money;
Freedom to Trade Internationally; Regulation), and
\item select subcomponents (pre-specified, high-salience, and non-overlapping with outcomes when needed).
\end{itemize}
(Area structure and components are defined in the EFW annual report. ([Fraser Institute][1]))

\vspace{0.4em}

\textbf{Macro reduced-form estimands.}\
For macro outcomes $Y$, define horizon-$h$ changes analogously:
\begin{equation}
\Delta Y_{ce}(h)
=
Y_{c,t_{ce}+h}
-
Y_{c,t_{ce}-1}.
\end{equation}

Primary reduced-form estimands:
\begin{equation}
\tau_{\Delta Y(h)}
=
\lim_{m \downarrow 0} \E\big[\Delta Y_{ce}(h)\given m_{ce}=m\big]
-
\lim_{m \uparrow 0} \E\big[\Delta Y_{ce}(h)\given m_{ce}=m\big].
\end{equation}

\vspace{0.4em}

\textbf{Tail / downside / crisis estimands (RD reduced form).}\
Define pre-specified tail metrics over the post-election horizon $[t_{ce},t_{ce}+h]$:
\begin{itemize}[leftmargin=1.5em]
\item \textbf{Growth downside:} $\Min_{s \in {0,\ldots,h}} g_{c,t_{ce}+s}$ where
$g_{c,t}$ is real GDP per capita growth.
\item \textbf{Inflation spike:} $\1{ \Max_{s \in {0,\ldots,h}} \pi_{c,t_{ce}+s} \ge \pi^\star}$,
with pre-specified thresholds $\pi^\star \in \{20\%,40\%\}$.
\item \textbf{Max drawdown:} peak-to-trough log decline in real GDP per capita within horizon:
\[
\mathrm{MDD}_{ce}(h)
=
\Max_{0 \le u \le v \le h}\ \big(\log y_{c,t_{ce}+u} - \log y_{c,t_{ce}+v}\big).
\]
\item \textbf{Crisis start:} $\1{\text{systemic banking crisis starts in }[t_{ce},t_{ce}+h]}$ from the
\LV\ database (1970--2011) where available; see data section for updates/fallbacks. ([SSRN][9])
\end{itemize}

For each tail metric $T_{ce}(h)$, define $\tau_{T(h)}$ as the RD discontinuity.

\vspace{0.8em}

%%%%%%%%%%%%%%%%%%%%%%%%%%%%%%%%%%%%%%%%%%%%%%%%%%%%%%%%%%%%%%%%%%%%%%%%%%%%%%%
\subsection*{1.2 Part 2 estimands}
\textit{(EFW $\rightarrow$ macro; required; explicit about localness and assumptions.)}

\vspace{0.4em}

\textbf{Objective:} identify the causal effect of election-induced EFW changes on macro outcomes,
allowing asymmetry, state dependence, and nonlinearities.

\vspace{0.4em}

\textbf{Baseline Part-2 estimand: local Wald (RD-IV) effect of EFW on macro at horizon $h$.}\
Define the fuzzy-RD / local-IV estimand:
\begin{equation}
\beta^{\IV}(h)
=
\frac{\tau_{\Delta Y(h)}}{\tau_{\Delta \EFW(h)}},
\label{eq:wald}
\end{equation}
where both numerator and denominator are RD discontinuities at $m=0$.

Interpretation:
\begin{itemize}[leftmargin=1.5em]
\item Under RD continuity + instrument relevance ($\tau_{\Delta \EFW(h)} \ne 0$),
plus an exclusion restriction that the close-election winner affects $\Delta Y(h)$ only through
its effect on $\Delta \EFW(h)$, the ratio identifies a local average causal effect of EFW on $Y$
for ``compliers'' at the margin of victory.
\item We will \textbf{not} present \eqref{eq:wald} as unconditional truth; we will present it as
(i) a local causal estimand under additional assumptions, and
(ii) a sensitivity/bounds object when exclusion is imperfect.
\end{itemize}

\vspace{0.4em}

\textbf{Defining ``increase vs decrease'' without post-treatment conditioning.}\
We must avoid selecting observations based on realized $\Delta\EFW$ sign.
Our solution is estimand design via \emph{instrumented piecewise effects} using
pre-treatment regime indicators and/or multiple instruments created by interacting $D_{ce}$ with
pre-treatment states.

We define:
\begin{equation}
\Delta \EFW^+_{ce}(h) = \max\{\Delta \EFW_{ce}(h),0\},
\qquad
\Delta \EFW^-_{ce}(h) = \min\{\Delta \EFW_{ce}(h),0\}.
\end{equation}

Then an asymmetric structural specification:
\begin{equation}
\Delta Y_{ce}(h)
=
\beta^+(h)\Delta \EFW^+_{ce}(h)
+
\beta^-(h)\Delta \EFW^-_{ce}(h)
+
u_{ce}(h),
\label{eq:asymstruct}
\end{equation}
where both components are treated as endogenous and instrumented using $D_{ce}$ and
\emph{interactions with predetermined states} (see Sections 6--8).

Key point:
\begin{itemize}[leftmargin=1.5em]
\item We do \textbf{not} condition on $\sign(\Delta \EFW)$ in sample selection.
\item We allow sign-dependent slopes by modeling $\Delta \EFW$ in two endogenous parts.
\item Identification requires at least two relevant instruments; we generate them by interacting
the close-election win with predetermined state variables that shift the sign/magnitude of EFW responses.
\end{itemize}

\vspace{0.4em}

\textbf{State-dependent causal parameters (CATE-IV / local-IV).}\
Let $S_{ce}$ be a predetermined state variable (or vector), measured before election:
baseline EFW, inflation history, democracy/institutions, openness, etc.
We target:
\begin{equation}
\beta^{\IV}(h;S=s)
=
\frac{\tau_{\Delta Y(h)\mid S=s}}{\tau_{\Delta \EFW(h)\mid S=s}},
\end{equation}
and interaction-IV analogues:
\begin{equation}
\Delta Y_{ce}(h)
=
\beta_0(h)\Delta \EFW_{ce}(h)
+
\beta_1(h)\big(\Delta \EFW_{ce}(h)\times S_{ce}\big)
+
u_{ce}(h),
\end{equation}
with instruments $(D_{ce}, D_{ce}\times S_{ce})$ under local RD conditions.

\vspace{0.4em}

\textbf{Nonlinear causal parameters.}\
We target threshold and magnitude nonlinearities:
\begin{itemize}[leftmargin=1.5em]
\item \textbf{Threshold in baseline level:} effects differ if pre-election $\EFW_{c,t_{ce}-1}$
is below/above $\kappa$ (pre-specified or cross-validated within RD-valid discipline).
\item \textbf{Magnitude nonlinearity:} marginal effect depends on $|\Delta \EFW|$ (small vs large reforms).
\item \textbf{Tail nonlinearity:} EFW affects probability/severity of extreme events more than mean outcomes.
\item \textbf{Component nonlinearity:} effects depend on which EFW area moves (e.g., Legal System vs Regulation).
\end{itemize}

All nonlinearities will be estimated via \emph{pre-specified instrumented piecewise models}
with multiplicity control (Section 8).

\vspace{0.8em}

%%%%%%%%%%%%%%%%%%%%%%%%%%%%%%%%%%%%%%%%%%%%%%%%%%%%%%%%%%%%%%%%%%%%%%%%%%%%%%%
\subsection*{1.3 Claims discipline}
\textit{(No overclaiming; explicit ladder of assumptions.)}

\vspace{0.4em}

\textbf{Tier 1 (minimal assumptions; primary claims).}\
Under standard close-election RD continuity assumptions, we can claim:
\begin{itemize}[leftmargin=1.5em]
\item $\tau_{\Delta \EFW(h)}$: local causal effect of more-market winner on EFW changes.
\item $\tau_{\Delta Y(h)}$ and $\tau_{T(h)}$: local causal effects of more-market winner on macro outcomes
and tail metrics.
\end{itemize}

\vspace{0.4em}

\textbf{Tier 2 (additional assumptions; cautious causal `EFW $\rightarrow$ macro'').}\\
Interpreting $\beta^{\IV}(h)$ as causal requires:
\begin{itemize}[leftmargin=1.5em]
\item \textbf{Relevance:} $\tau_{\Delta \EFW(h)} \neq 0$ within the estimation window.
\item \textbf{Exclusion:} the winner affects $\Delta Y(h)$ only through $\Delta \EFW(h)$
(or through modeled EFW components), locally at the cutoff.
\item \textbf{Monotonicity (if using LATE language):} no `defiers'' in terms of EFW response to
electing the more-market competitor (must be argued and partially diagnosed).
\item \textbf{SUTVA/local stability:} no interference across countries/elections that violates the local design.
\end{itemize}

We will treat these as strong assumptions, and we will \textbf{not} hide them.

\vspace{0.4em}

\textbf{Tier 3 (exploratory but disciplined).}\
Nonlinearities/asymmetries/state dependence in $\beta$ will be:
\begin{itemize}[leftmargin=1.5em]
\item pre-specified where possible,
\item corrected for multiple testing,
\item reported with weak-IV robust inference,
\item accompanied by sensitivity analysis for exclusion and measurement error.
\end{itemize}

\vspace{1.0em}

%%%%%%%%%%%%%%%%%%%%%%%%%%%%%%%%%%%%%%%%%%%%%%%%%%%%%%%%%%%%%%%%%%%%%%%%%%%%%%%
\section*{2) Identification Strategy --- Part 1}
\textit{(Improved RD; flagship spine.)}

\vspace{0.4em}

%%%%%%%%%%%%%%%%%%%%%%%%%%%%%%%%%%%%%%%%%%%%%%%%%%%%%%%%%%%%%%%%%%%%%%%%%%%%%%%
\subsection*{2.1 Unit, timing, running variable, treatment, orientation}

\vspace{0.3em}

\textbf{Unit of analysis.}\
Default unit: a national election that determines the head of government or governing coalition,
indexed by $(c,e)$ with election year $t_{ce}$.

\vspace{0.3em}

\textbf{Baseline timing.}\
Baseline year is $t_{ce}-1$ (pre-election).
We map outcomes at horizons $h \in {0,1,2,3,4,5}$ (annual) and, where needed,
$h \in {5}$ as a quinquennial endpoint for pre-2000 EFW coverage.

\vspace{0.3em}

\textbf{Running variable definitions (realistic by electoral system).}\
We define $m_{ce}$ as the \emph{signed} margin between more-market and less-market competitors
in the decisive contest that assigns executive power:

\begin{itemize}[leftmargin=1.5em]
\item \textbf{Presidential two-round systems:}
$m_{ce}$ is the vote-share difference in the \emph{final decisive round} between the two finalists
mapped as more-market vs less-market.
\item \textbf{Presidential one-round plurality:}
$m_{ce}$ is the vote-share difference between the top two candidates, with the winner mapped as
more-market or less-market relative to the runner-up.
\item \textbf{Parliamentary plurality/PR:}
default $m_{ce}$ is the seat-share margin between the two largest competing government-forming
blocs, defined via pre-election coalition declarations where available, or via the largest party
that forms government vs the largest party that does not.
We restrict to cases where government formation is not ambiguous (see restrictions below).
\item \textbf{Runoff / coalition ambiguity:}
if the decisive ``win'' is not pinned to a clear margin (e.g., protracted coalition bargaining),
the observation is excluded from the baseline sample and used only in a robustness appendix
with alternative running variables (e.g., coalition seat margin around 50\%).
\end{itemize}

\vspace{0.3em}

\textbf{Treatment indicator.}\
$D_{ce} = \1{m_{ce} > 0}$ indicates a more-market winner.
The complement is a less-market winner.

\vspace{0.3em}

\textbf{Orientation rule: more-market vs less-market competitor.}\
We define an ideology/market-orientation score $I_{p,t}$ for party/candidate $p$ at time $t$.
Higher $I$ means more market-oriented.

Baseline orientation algorithm (pre-specified):
\begin{enumerate}[leftmargin=1.8em]
\item \textbf{Primary source: V-Dem V-Party party economic positions} where coverage exists,
mapped to election year. ([datafinder.qog.gu.se][12])
\item \textbf{Secondary source: World Bank DPI ideology variables} (government and opposition orientation)
for broader country coverage. ([Inter-American Development Bank][4])
\item \textbf{Tertiary source: Manifesto Project (MARPOR) left-right / economic scales}
for countries with coded manifestos. ([CRAN][6])
\item \textbf{Fallback: hand-coded ``market orientation'' using transparent rules}
restricted to a small set of high-stakes elections, with ex-ante coding and blind double coding.
(Used only if the above fail and reported transparently.)
\end{enumerate}

Orientation in election $(c,e)$:
\begin{itemize}[leftmargin=1.5em]
\item Identify the top two competitors in decisive contest (candidate or party).
\item Assign each competitor a pre-election $I$ score.
\item Label the competitor with higher $I$ as `more-market'' and the other as `less-market''.
\item If $|I_A - I_B| < \delta_I$ (pre-specified minimal ideological separation), drop from baseline
to avoid classification noise; include in robustness with continuous $I$.
\end{itemize}

\vspace{0.3em}

\textbf{Sample restrictions (baseline).}\
To preserve RD credibility and comparability:
\begin{itemize}[leftmargin=1.5em]
\item Keep elections with credible competitiveness and reliable margins (official/validated sources).
\item Exclude elections with documented irregularities severe enough to plausibly affect the margin.
(Operationalized via external democracy/election-quality indicators; see state variables section.)
\item Exclude parliamentary elections with ambiguous government formation (unless a clean
pre-election coalition exists).
\item Exclude elections too close to EFW missingness or where EFW data are unavailable for required horizon.
\end{itemize}

\vspace{0.8em}

%%%%%%%%%%%%%%%%%%%%%%%%%%%%%%%%%%%%%%%%%%%%%%%%%%%%%%%%%%%%%%%%%%%%%%%%%%%%%%%
\subsection*{2.2 RD improvements}
\textit{(Propose and choose a default best-practice RD design.)}

\vspace{0.4em}

\textbf{Core challenge:} cross-country close-election RD involves
(i) heterogeneous electoral systems,
(ii) repeated elections per country,
(iii) serial correlation in outcomes,
(iv) discrete/heaped running variables,
(v) multi-horizon outcome construction, and
(vi) potential manipulation/fraud in a subset of contexts.

\vspace{0.4em}

\textbf{Default estimation strategy: robust bias-corrected local polynomial RD.}\
We adopt modern RD inference with robust bias correction and data-driven bandwidth selection. ([rdpackages.github.io][10])\
We implement:
\begin{itemize}[leftmargin=1.5em]
\item local linear (order 1) as baseline point estimator with bias correction using local quadratic,
\item MSE-optimal bandwidths for the bias-corrected estimator,
\item robust standard errors clustered by country (primary) and, where needed, two-way clustering
(country $\times$ calendar year) in pooled-year specifications.
\end{itemize}

\vspace{0.4em}

\textbf{Complementary robustness: local randomization inference.}\
We treat observations within a narrow window $|m_{ce}| \le r$ as a locally randomized experiment,
using finite-sample randomization inference and window selection based on covariate balance. ([rdpackages.github.io][13])\
This is not the default because window choice is delicate, but it is a strong robustness check when
sample sizes are modest near the cutoff.

\vspace{0.4em}

\textbf{Bandwidth strategy.}\
We pre-specify:
\begin{itemize}[leftmargin=1.5em]
\item Main results: data-driven MSE-optimal bandwidths under \CCT\ robust inference.
\item Sensitivity: half-bandwidth and double-bandwidth.
\item ``Honest'' reporting: show bandwidth-sensitivity plots for $\tau$ across a grid.
\item Local randomization: choose symmetric windows via covariate-balance criteria and report p-values
over a range of $r$.
\end{itemize}

\vspace{0.4em}

\textbf{Dependence handling.}\
Multiple elections per country induce within-country dependence and potentially overlapping horizons.
We address this at design and inference levels:

\begin{itemize}[leftmargin=1.5em]
\item \textbf{Primary inference:} country-clustered standard errors / wild cluster bootstrap.
\item \textbf{Alternative:} collapse to one election per country per fixed calendar block (e.g., 4-year bins)
as a robustness check.
\item \textbf{Alternative:} treat each country as a ``clustered time series'' and implement block bootstrap
over election episodes within country (appendix).
\end{itemize}

\vspace{0.4em}

\textbf{Overlapping elections contamination (treatment switching).}\
Problem: If a subsequent election occurs within horizon $h$, the initially assigned winner may be replaced,
confounding interpretation of $\Delta Y_{ce}(h)$ and $\Delta \EFW_{ce}(h)$.

We adopt a hierarchy:

\begin{enumerate}[leftmargin=1.8em]
\item \textbf{Primary rule (term-aligned outcomes):} define $h$ relative to the expected term length
or until next election, whichever comes first; estimate effects on
\emph{within-term averages} and \emph{within-term cumulative changes}.
\item \textbf{Secondary rule (clean windows):} restrict to elections with no subsequent election within $h$
for selected horizons (e.g., $h=2$ and $h=4$), reported as a robustness sample.
\item \textbf{Intention-to-treat (ITT) interpretation:} for fixed calendar horizons regardless of switching,
interpret as ITT of initial close-election win, acknowledging dilution.
\end{enumerate}

We will not condition on post-election events as controls in primary specifications (to avoid
post-treatment conditioning).

\vspace{0.4em}

\textbf{Falsification suite tailored to elections.}\
We pre-specify falsifications:
\begin{itemize}[leftmargin=1.5em]
\item pre-election EFW trends (placebo outcomes at $t_{ce}-k$ for $k\ge 1$),
\item pre-election macro outcomes and tail metrics,
\item predetermined covariates: GDP per capita level, population, democracy index, trade openness,
historical inflation mean (all pre-election),
\item placebo cutoffs in the running variable (e.g., $\pm 2\%$ margins),
\item donut RD excluding ultra-close margins potentially most subject to manipulation or recount idiosyncrasies.
\end{itemize}

\vspace{0.8em}

%%%%%%%%%%%%%%%%%%%%%%%%%%%%%%%%%%%%%%%%%%%%%%%%%%%%%%%%%%%%%%%%%%%%%%%%%%%%%%%
\subsection*{2.3 Diagnostics}
\textit{(Manipulation, balance, specification, and external validity checks.)}

\vspace{0.4em}

\textbf{Manipulation/density.}\
We test for discontinuities in the density of $m_{ce}$ at 0 (with methods appropriate to discrete
running variables and heaping).
We interpret evidence cautiously in multi-country settings where rounding rules differ.

\vspace{0.4em}

\textbf{Covariate balance.}\
We estimate RD discontinuities in predetermined covariates:
\begin{itemize}[leftmargin=1.5em]
\item baseline $\EFW_{c,t_{ce}-1}$ and pre-trends,
\item baseline log GDP per capita,
\item baseline inflation average (past 3--5 years),
\item baseline democracy/institutions indices,
\item baseline trade openness and external vulnerability.
\end{itemize}
We require joint balance (multivariate test) within the chosen bandwidth/window.

\vspace{0.4em}

\textbf{Donut RD.}\
Exclude a small neighborhood $|m_{ce}| < \eta$ to address recount/manipulation concerns;
report $\eta$ sensitivity.

\vspace{0.4em}

\textbf{Placebo cutoffs and outcomes.}\
Estimate discontinuities at fake cutoffs and on outcomes measured prior to election.

\vspace{0.4em}

\textbf{Specification sensitivity.}\
Compare:
\begin{itemize}[leftmargin=1.5em]
\item local linear vs local quadratic,
\item triangular vs uniform kernels (where implemented),
\item robust bias correction vs conventional,
\item alternative bandwidth choices.
\end{itemize}

\vspace{0.4em}

\textbf{Heterogeneous electoral systems.}\
Report separate results by election type (presidential vs parliamentary) as a pre-specified heterogeneity
analysis (not fishing), with pooled main estimate as primary.

\vspace{1.0em}

%%%%%%%%%%%%%%%%%%%%%%%%%%%%%%%%%%%%%%%%%%%%%%%%%%%%%%%%%%%%%%%%%%%%%%%%%%%%%%%
\section*{3) Data Sources}
\textit{(No cleaning/ETL steps; sources + feasibility only; ranked options and fallbacks.)}

\vspace{0.4em}

%%%%%%%%%%%%%%%%%%%%%%%%%%%%%%%%%%%%%%%%%%%%%%%%%%%%%%%%%%%%%%%%%%%%%%%%%%%%%%%
\subsection*{3.1 EFW}

\textbf{Primary dataset: Economic Freedom of the World (EFW) dataset (Fraser Institute).}\
\begin{itemize}[leftmargin=1.5em]
\item \textbf{Provider / link:} Fraser Institute; `Economic Freedom of the World: 2025 Annual Report''
and associated dataset download. :contentReference[oaicite:22]{index=22}
\item \textbf{Coverage:} up to 165 jurisdictions as far back as 1970. :contentReference[oaicite:23]{index=23}
\item \textbf{Frequency:} mixed historical frequency (quinquennial early years) and annual in later years;
we treat pre-2000 primarily via quinquennial endpoints and 2000+ via annual horizons
(reported explicitly as distinct samples).
\item \textbf{Unit:} jurisdiction-year.
\item \textbf{Content:} overall index + 5 areas + 45 components/subcomponents. :contentReference[oaicite:24]{index=24}
\item \textbf{Access method:} free download from Fraser Institute dataset page; the report recommends using
the most recent dataset as `most up-to-date and accurate.'' ([Fraser Institute][1])
\item \textbf{Licensing/constraints:} academic use is standard; citation requested in the report
(the report provides a formal citation). ([Fraser Institute][1])
\item \textbf{Limitations:}
\begin{itemize}[leftmargin=1.5em]
\item EFW is a composite index; changes reflect both policy changes and data revisions.
\item Some components may embed macro outcomes (e.g., inflation in ``Sound Money''), creating mechanical
overlap with certain macro dependent variables.
\item Data sources and variable definitions evolve over time; comparability across long horizons requires
using a consistent vintage and noting breaks. ([Fraser Institute][1])
\end{itemize}
\end{itemize}

\textbf{Ranked fallbacks / complements.}\
\begin{itemize}[leftmargin=1.5em]
\item \textbf{Fallback 1:} Use EFW area-specific indices and exclude mechanically overlapping components
(e.g., exclude Sound Money when inflation is an outcome).
\item \textbf{Fallback 2:} Use EFW ranks/percentiles as robustness if level comparability is questioned.
\item \textbf{Complement:} Treat EFW as multiple latent dimensions (5 areas) rather than a single scalar;
estimate effects component-by-component.
\end{itemize}

\vspace{0.8em}

%%%%%%%%%%%%%%%%%%%%%%%%%%%%%%%%%%%%%%%%%%%%%%%%%%%%%%%%%%%%%%%%%%%%%%%%%%%%%%%
\subsection*{3.2 Elections + margins (at least 3 sources, ranked)}

\textbf{Source 1 (global, legislative; high coverage): Constituency-Level Elections Archive (CLEA).}\
\begin{itemize}[leftmargin=1.5em]
\item \textbf{Provider / link:} electiondataarchive.org (CLEA project). ([electiondataarchive.org][2])
\item \textbf{Coverage:} global constituency-level results for lower house legislative elections;
latest update noted as Oct 15, 2025 with 2,181 elections in 183 countries and large-scale returns. ([electiondataarchive.org][2])
\item \textbf{Unit:} constituency-level election returns; can be aggregated to national vote/seat totals.
\item \textbf{Access method:} downloadable files (archive), plus documentation; re3data notes it is a
public good ``publicly available at no cost'' though some access may require registration. ([re3data.org][14])
\item \textbf{Licensing/constraints:} heterogeneity across underlying sources; attribution/citation required;
some country returns may have specific restrictions; we will comply with repository terms.
\item \textbf{Limitations:}
\begin{itemize}[leftmargin=1.5em]
\item Primarily legislative elections; presidential contests may need separate sources.
\item Coalition formation is not directly observed; mapping to executive requires additional political data.
\item Differences in electoral rules and district structures require careful margin definitions.
\end{itemize}
\end{itemize}

\textbf{Source 2 (EU/OECD parliamentary; clean executive mapping): ParlGov.}\
\begin{itemize}[leftmargin=1.5em]
\item \textbf{Provider / link:} ParlGov project and data releases. ([parlgov.org][3])
\item \textbf{Coverage:} EU and OECD democracies 1900--2023; includes parties, elections, cabinets. ([parlgov.org][3])
\item \textbf{Unit:} country-election, party vote and seat shares; cabinet/government spells.
\item \textbf{Access method:} CSV download (e.g., GitHub releases) and Harvard Dataverse. ([parlgov.org][15])
\item \textbf{Licensing/constraints:} generally open academic infrastructure; confirm per release.
\item \textbf{Limitations:}
\begin{itemize}[leftmargin=1.5em]
\item Restricted to EU/OECD democracies (not global).
\item Executive ``winner'' in parliamentary systems may be mediated by coalition bargaining.
\end{itemize}
\end{itemize}

\textbf{Source 3 (broad, political institutions; ideology linkage): Database of Political Institutions (DPI).}\
\begin{itemize}[leftmargin=1.5em]
\item \textbf{Provider / link:} World Bank DPI (e.g., DPI2020). ([Inter-American Development Bank][4])
\item \textbf{Coverage:} 1975--2020, 180 countries, political institutions and checks and balances;
includes party/ideology measures. ([Inter-American Development Bank][4])
\item \textbf{Unit:} country-year (political institutions) with variables that can be matched to election years.
\item \textbf{Access method:} World Bank data catalog / downloads.
\item \textbf{Licensing/constraints:} World Bank terms; citation requirement; check access terms in catalog.
\item \textbf{Limitations:}
\begin{itemize}[leftmargin=1.5em]
\item Not a detailed election-results database; margins may be coarse or absent for some contests.
\item Ideology measures are categorical/coarse in many cases, useful as fallback not primary.
\end{itemize}
\end{itemize}

\textbf{Ranked fallback options for elections/margins (if needed).}\
\begin{itemize}[leftmargin=1.5em]
\item \textbf{Fallback A:} Restrict baseline sample to contexts with clean margins (e.g., presidential
top-two final round, or parliamentary single-party majority) even if it reduces coverage.
\item \textbf{Fallback B:} Use ParlGov for EU/OECD and treat rest-of-world as separate paper/appendix.
\item \textbf{Fallback C:} Use multiple sources and reconcile margins by ``two-out-of-three'' agreement;
flag discrepancies and run robustness excluding discrepant cases.
\end{itemize}

\vspace{0.8em}

%%%%%%%%%%%%%%%%%%%%%%%%%%%%%%%%%%%%%%%%%%%%%%%%%%%%%%%%%%%%%%%%%%%%%%%%%%%%%%%
\subsection*{3.3 ``More-market vs less-market'' scoring (3 strategies ranked)}

\textbf{Strategy 1 (preferred where available): V-Dem V-Party expert-coded party positions.}\
\begin{itemize}[leftmargin=1.5em]
\item \textbf{Provider / link:} V-Dem V-Party dataset and documentation; accessible via V-Dem pages. ([datafinder.qog.gu.se][12])
\item \textbf{Coverage:} broad cross-national coverage with party positions; QoG summary indicates
coverage 1970--2019 and inclusion of parties with $>5\%$ vote shares, across 169 countries. ([datafinder.qog.gu.se][12])
\item \textbf{Unit:} party-election or party-year (depending on variable); expert-coded positions.
\item \textbf{Access method:} dataset download (formats STATA/CSV etc) from V-Dem; requires matching to
election year and party identifiers. ([v-dem.net][5])
\item \textbf{What it measures:} party economic left-right and related dimensions, usable to define
``more-market'' ordering within an election.
\item \textbf{Constraints/limitations:}
\begin{itemize}[leftmargin=1.5em]
\item Expert-coded measures have uncertainty and may not cover very small parties.
\item Party identifiers must be reconciled with election data; feasible but nontrivial (no ETL details here).
\end{itemize}
\end{itemize}

\textbf{Strategy 2 (global fallback): DPI ideology categories.}\
\begin{itemize}[leftmargin=1.5em]
\item \textbf{Provider / link:} World Bank DPI2020. ([Inter-American Development Bank][4])
\item \textbf{Coverage:} 1975--2020, 180 countries; includes ideology measures for key actors. ([Inter-American Development Bank][4])
\item \textbf{Unit:} country-year.
\item \textbf{What it measures:} coarse left/center/right or similar categorizations for executive/parties.
\item \textbf{Use case:} define ``more-market'' as right/market-liberal vs left/interventionist, with
explicit acknowledgment of coarseness.
\item \textbf{Limitations:} ideology coding may not track market reform stance in all regions/time periods;
we treat as robustness/coverage extender.
\end{itemize}

\textbf{Strategy 3 (high validity where available): Manifesto Project (MARPOR) text-based scales.}\
\begin{itemize}[leftmargin=1.5em]
\item \textbf{Provider / link:} Manifesto Project; access via manifestoR and API key. ([CRAN][6])
\item \textbf{Coverage:} strong for many European and selected non-European democracies with coded manifestos;
time coverage varies by country.
\item \textbf{Unit:} party-election manifesto coding.
\item \textbf{Access method:} requires registration/API key; manifestoR provides retrieval and versioning. ([CRAN][6])
\item \textbf{What it measures:} issue-category shares and derived left-right/economic scales.
\item \textbf{Limitations:}
\begin{itemize}[leftmargin=1.5em]
\item Limited global coverage; not universal across developing democracies.
\item Manifesto coding captures stated priorities, not necessarily implemented policy.
\end{itemize}
\end{itemize}

\textbf{Feasibility ranking (orientation).}\
\begin{itemize}[leftmargin=1.5em]
\item Best validity: V-Party and MARPOR where available, triangulated.
\item Best coverage: DPI ideology where finer sources missing.
\item Baseline: use V-Party for main sample where possible; expand via DPI for breadth; report
MARPOR-based subsample as ``high-precision'' robustness.
\end{itemize}

\vspace{0.8em}

%%%%%%%%%%%%%%%%%%%%%%%%%%%%%%%%%%%%%%%%%%%%%%%%%%%%%%%%%%%%%%%%%%%%%%%%%%%%%%%
\subsection*{3.4 Macro outcomes + tails + crises/drawdowns}

\textbf{Primary macro source 1: World Development Indicators (WDI) via World Bank API.}\
\begin{itemize}[leftmargin=1.5em]
\item \textbf{Provider / link:} World Bank WDI API documentation and licensing. ([datahelpdesk.worldbank.org][7])
\item \textbf{Coverage:} broad country coverage; annual macro indicators (varies by series and country).
\item \textbf{Unit:} country-year.
\item \textbf{Access method:} API without keys; standard JSON/XML retrieval; convenient for replication. ([datahelpdesk.worldbank.org][7])
\item \textbf{License:} WDI data are available under CC BY 4.0 per World Bank licensing page. ([datacatalog.worldbank.org][16])
\item \textbf{Core indicators (examples):}
\begin{itemize}[leftmargin=1.5em]
\item Real GDP growth (annual \%), GDP per capita, CPI inflation, gross fixed capital formation (\% GDP),
trade openness (\% GDP), population.
\end{itemize}
\item \textbf{Limitations:}
\begin{itemize}[leftmargin=1.5em]
\item Measurement comparability and revisions across time.
\item Missingness for fragile states and high-inflation episodes; we pre-specify missing-data handling in analysis.
\end{itemize}
\end{itemize}

\textbf{Primary macro source 2: Penn World Table (PWT).}\
\begin{itemize}[leftmargin=1.5em]
\item \textbf{Provider / link:} Groningen Growth and Development Centre (GGDC) PWT releases.
PWT 10.0 provides Excel/Stata downloads and is CC BY 4.0. ([University of Groningen][8])
\item \textbf{Coverage:} PWT 10.0 covers 183 countries, 1950--2019 (per release page). ([University of Groningen][8])
\item \textbf{Update:} PWT 11.0 released Oct 7, 2025 and covers 185 countries through 2023. ([University of Groningen][17])
\item \textbf{Unit:} country-year.
\item \textbf{Access method:} downloadable datasets; (optional) replication packages provided. ([University of Groningen][8])
\item \textbf{License:} CC BY 4.0 (PWT 10.0 page). ([University of Groningen][8])
\item \textbf{Use case:} real GDP per capita levels and growth, investment/capital measures, productivity.
\item \textbf{Limitations:} PPP concepts differ from national accounts; revisions across versions;
we fix a vintage for main results and show robustness to alternative series.
\end{itemize}

\textbf{Tail and drawdown construction (from macro series).}\
We construct drawdowns and tail metrics directly from annual GDP per capita and inflation series,
avoiding dependence on proprietary crisis chronologies when possible.

\textbf{Crisis dataset option (binary tail).}\
\begin{itemize}[leftmargin=1.5em]
\item \textbf{Laeven--Valencia Systemic Banking Crises Database (update).}
Coverage includes systemic banking, currency, and sovereign debt crises 1970--2011. ([SSRN][9])
\item \textbf{Limitations:} ends in 2011 per the cited update; for post-2011, we will (i) rely on
macro-constructed crisis proxies (inflation spikes, large drawdowns), and/or (ii) use later public
updates if accessible (to be documented in replication archive).
\end{itemize}

\vspace{0.8em}

%%%%%%%%%%%%%%%%%%%%%%%%%%%%%%%%%%%%%%%%%%%%%%%%%%%%%%%%%%%%%%%%%%%%%%%%%%%%%%%
\subsection*{3.5 Pre-treatment state variables (for state dependence)}

\textbf{Primary institutional/democracy source: V-Dem country-year dataset.}\
\begin{itemize}[leftmargin=1.5em]
\item \textbf{Provider / link:} V-Dem Dataset download page.
Version 15 published March 2025; free downloads include codebook and cautionary notes. ([v-dem.net][18])
\item \textbf{Coverage:} extensive country-year democracy and institutional indicators.
\item \textbf{Unit:} country-year (core and full datasets).
\item \textbf{Use case:} pre-election state variables:
democracy indices, constraints on executive, rule of law proxies, corruption indicators,
civil liberties, and election-quality measures.
\item \textbf{Constraints:} V-Dem warns against comparing absolute index values across versions;
we fix version 15 (March 2025) as baseline. ([v-dem.net][18])
\end{itemize}

\textbf{State capacity / governance (ranked options).}\
\begin{itemize}[leftmargin=1.5em]
\item \textbf{Option A (broad, common):} governance effectiveness and regulatory quality indicators
(e.g., Worldwide Governance Indicators) as predetermined states.
\item \textbf{Option B (fiscal capacity):} tax revenue-to-GDP (WDI/IMF sources) as capacity proxy.
\item \textbf{Option C (bureaucratic quality):} V-Dem administrative capacity measures (where available).
\end{itemize}

\textbf{Openness and vulnerability (predetermined states).}\
\begin{itemize}[leftmargin=1.5em]
\item Trade openness (exports+imports)/GDP from WDI.
\item Financial openness proxies (where available) and external debt indicators (WDI/other).
\item Inflation history: average and volatility of inflation over $t_{ce}-5$ to $t_{ce}-1$ (WDI).
\item Baseline income: log GDP per capita at $t_{ce}-1$ (WDI/PWT).
\end{itemize}

\vspace{1.0em}

%%%%%%%%%%%%%%%%%%%%%%%%%%%%%%%%%%%%%%%%%%%%%%%%%%%%%%%%%%%%%%%%%%%%%%%%%%%%%%%
\section*{4) Empirical Specifications}
\textit{(Equations + what is primary; horizons; tails defined precisely.)}

\vspace{0.4em}

\textbf{Part 1: RD reduced-form and first-stage equations (primary).}\
For each horizon $h$ and outcome $O_{ce}(h)$, estimate:
\begin{equation}
O_{ce}(h)
=
\alpha_h
+
\tau_h D_{ce}
+
f^-_h(m_{ce})\cdot \1{m_{ce}<0}
+
f^+_h(m_{ce})\cdot \1{m_{ce}\ge 0}
+
\eps_{ce}(h),
\label{eq:rd}
\end{equation}
where $f^-_h$ and $f^+_h$ are local polynomials (default local linear) estimated within bandwidth
$\bw_h$ selected by robust MSE-optimal criteria, with bias correction. ([rdpackages.github.io][10])

Primary outcomes in Part 1:
\begin{itemize}[leftmargin=1.5em]
\item $O_{ce}(h)=\Delta \EFW_{ce}(h)$ (overall and components).
\item $O_{ce}(h)=\Delta Y_{ce}(h)$ for macro outcomes.
\item $O_{ce}(h)=T_{ce}(h)$ for tail metrics (binary or continuous).
\end{itemize}

\vspace{0.4em}

\textbf{Timing/event-time mapping and horizons.}\
We pre-specify horizons $h$ as:
\begin{itemize}[leftmargin=1.5em]
\item Annual: $h \in {0,1,2,3,4,5}$ (where data coverage permits).
\item Quinquennial regime: for earlier EFW frequency, focus on $h=5$ endpoints using
$\EFW_{c,t}$ at quinquennial years (sample-specific).
\item Term-aligned: define $h$ as ``years into term'' up to min(term length, 5) for comparability.
\end{itemize}

We explicitly report which sample uses which horizon mapping.

\vspace{0.4em}

\textbf{Tail outcomes (precise definitions).}\
Let $y_{c,t}$ be real GDP per capita.
Let $g_{c,t} = 100\cdot(\log y_{c,t} - \log y_{c,t-1})$.
Let $\pi_{c,t}$ be CPI inflation rate (annual \%).

Define for each election $(c,e)$ and horizon $h$:

\begin{align}
\mathrm{WorstGrowth}_{ce}(h)
&=
\Min_{s=0,\ldots,h}\ g_{c,t_{ce}+s},
\\
\mathrm{InflSpike}_{ce}(h;\pi^\star)
&=
\1\Big\{\Max_{s=0,\ldots,h}\ \pi_{c,t_{ce}+s} \ge \pi^\star\Big\},
\qquad
\pi^\star \in \{20,40\},
\\
\mathrm{MDD}_{ce}(h)
&=
\Max_{0 \le u \le v \le h}\ \big(\log y_{c,t_{ce}+u} - \log y_{c,t_{ce}+v}\big),
\\
\mathrm{Crisis}_{ce}(h)
&=
\1\big\{\exists t \in [t_{ce},t_{ce}+h] \text{ s.t. systemic banking crisis starts at } t\big\},
\end{align}
where $\mathrm{Crisis}_{ce}(h)$ is measured from the \LV\ crisis chronology where available. ([SSRN][9])

\vspace{0.4em}

\textbf{Primary inference reporting.}\
We report:
\begin{itemize}[leftmargin=1.5em]
\item RD point estimate $\hat{\tau}_h$ and robust bias-corrected 95\% CI.
\item Effective sample sizes on each side (within bandwidth).
\item Clustered SE at country level and, as robustness, wild cluster bootstrap.
\item Sensitivity to bandwidth multiples and polynomial order.
\end{itemize}

\vspace{0.4em}

\textbf{Avoiding tautology for overlapping variables.}\
Because EFW includes inflation-related measures in the Sound Money area, we:
\begin{itemize}[leftmargin=1.5em]
\item do not interpret EFW $\rightarrow$ inflation results using overall EFW without adjustment,
\item instead use ``EFW excluding Sound Money'' or use non-overlapping EFW areas when inflation is the
dependent variable,
\item and report component-wise effects to show which institutional dimensions drive macro results. ([Fraser Institute][1])
\end{itemize}

\vspace{1.0em}

%%%%%%%%%%%%%%%%%%%%%%%%%%%%%%%%%%%%%%%%%%%%%%%%%%%%%%%%%%%%%%%%%%%%%%%%%%%%%%%
\section*{PART 2 (REQUIRED) --- EFW $\rightarrow$ Macro}
\textit{(Asymmetry + state dependence + nonlinearities; identification-first.)}

\vspace{0.6em}

%%%%%%%%%%%%%%%%%%%%%%%%%%%%%%%%%%%%%%%%%%%%%%%%%%%%%%%%%%%%%%%%%%%%%%%%%%%%%%%
\section*{5) Part 2 Identification}
\textit{(How to identify EFW $\rightarrow$ macro from close elections; $\ge 3$ strategies;
assumptions; failure modes; robustness/bounds; choose default.)}

\vspace{0.4em}

\textbf{Part-2 problem statement.}\
Part 1 estimates $\tau_{\Delta \EFW(h)}$ and $\tau_{\Delta Y(h)}$ at close-election cutoff.
Part 2 seeks a causal effect of EFW on macro outcomes, not merely winner effects.

Core difficulty:
\begin{itemize}[leftmargin=1.5em]
\item Winner $D_{ce}$ can affect macro outcomes through channels beyond EFW (policy bundles, expectations,
geopolitical actions, etc.).
\item Therefore, `$D \rightarrow Y$'' is not automatically interpretable as `$\EFW \rightarrow Y$''.
\end{itemize}

We propose three identification strategies.

\vspace{0.6em}

%%%%%%%%%%%%%%%%%%%%%%%%%%%%%%%%%%%%%%%%%%%%%%%%%%%%%%%%%%%%%%%%%%%%%%%%%%%%%%%
\subsection*{5.1 Candidate Strategy 1: Local Wald RD-IV (ratio of discontinuities)}

\textbf{Estimand.}\
For horizon $h$:
\[
\beta^{\IV}(h)=\tau_{\Delta Y(h)}/\tau_{\Delta \EFW(h)}.
\]

\textbf{Assumptions.}
\begin{itemize}[leftmargin=1.5em]
\item RD continuity: potential outcomes and potential EFW paths are continuous in $m$ at 0.
\item Relevance: $\tau_{\Delta \EFW(h)}\neq 0$ (first stage at horizon $h$).
\item Exclusion: $D$ affects $\Delta Y(h)$ only through $\Delta \EFW(h)$ (or through specified EFW components).
\item Monotonicity (optional for LATE): no units whose EFW responds in opposite direction to $D$.
\item Stable units: no interference; local SUTVA.
\end{itemize}

\textbf{Failure modes (skeptical referee).}
\begin{itemize}[leftmargin=1.5em]
\item Exclusion is the main concern: elected officials affect macro through non-EFW channels
(e.g., discrete fiscal stimulus, commodity policy, conflict, appointments, central bank independence)
that may not be captured by EFW.
\item ``Bad control'' risk: if we try to control for post-election policies to satisfy exclusion,
we introduce post-treatment bias.
\item Weak first stage for some horizons or subsets reduces credibility (weak-IV distortions).
\end{itemize}

\textbf{Robustness/sensitivity/bounds (at least 2).}
\begin{enumerate}[leftmargin=1.8em]
\item \textbf{Plausibly exogenous IV sensitivity (Conley-type):}
allow $D$ to have a direct effect on $Y$ bounded within $[\ubarDelta,\barDelta]$ and compute
identified set for $\beta^{\IV}(h)$.
\item \textbf{Component-exclusion tightening:}
use multiple endogenous EFW areas and interpret $\beta$ as effect of specific institutional channels;
test whether results hinge on areas likely to embed direct macro outcomes (e.g., Sound Money vs inflation).
\item \textbf{Placebo ``direct effects'':}
estimate RD effects on outcomes that should not respond to policy within short horizons (e.g.,
pre-election outcomes, predetermined demographics) to calibrate plausible magnitude of direct effects.
\item \textbf{Weak-IV robust inference:} Anderson--Rubin style tests and confidence sets for $\beta$.
\end{enumerate}

\vspace{0.6em}

%%%%%%%%%%%%%%%%%%%%%%%%%%%%%%%%%%%%%%%%%%%%%%%%%%%%%%%%%%%%%%%%%%%%%%%%%%%%%%%
\subsection*{5.2 Candidate Strategy 2: 2SLS within bandwidth (fuzzy RD with controls)}

\textbf{Estimand.}\
Within a chosen bandwidth $\abs{m_{ce}}\le \bw$, estimate:
\begin{align}
\Delta \EFW_{ce}(h)
&=
\pi_0(h)+\pi_1(h)D_{ce}+g(m_{ce})+\nu_{ce}(h),
\\
\Delta Y_{ce}(h)
&=
\alpha_0(h)+\beta(h)\Delta \EFW_{ce}(h)+g(m_{ce})+u_{ce}(h),
\end{align}
instrumenting $\Delta \EFW_{ce}(h)$ with $D_{ce}$ and using the same RD control function $g(\cdot)$.

\textbf{Assumptions.}
Same as Strategy 1, but with additional parametric structure:
\begin{itemize}[leftmargin=1.5em]
\item Correct specification of local polynomial $g(m)$ within bandwidth (mitigated by local approach).
\item Exclusion as above, possibly conditional on predetermined covariates (covariate-adjusted RD).
\end{itemize}

\textbf{Failure modes (skeptical referee).}
\begin{itemize}[leftmargin=1.5em]
\item Results may depend on bandwidth choice and polynomial order.
\item If we add many controls, we risk finite-sample instability and p-hacking concerns.
\end{itemize}

\textbf{Robustness/sensitivity/bounds.}
\begin{enumerate}[leftmargin=1.8em]
\item Use robust bias correction in both stages (fuzzy RD implementation) and report sensitivity.
\item Use local randomization IV as an alternative (Strategy 3 below).
\item Apply weak-IV robust tests and report first-stage F-statistics (cluster-robust analogues).
\end{enumerate}

\vspace{0.6em}

%%%%%%%%%%%%%%%%%%%%%%%%%%%%%%%%%%%%%%%%%%%%%%%%%%%%%%%%%%%%%%%%%%%%%%%%%%%%%%%
\subsection*{5.3 Candidate Strategy 3: Local randomization + IV (finite-sample encouragement design)}

\textbf{Estimand.}\
Within window $\abs{m_{ce}}\le r$, treat $D_{ce}$ as randomized with probability near 0.5.
Then estimate an IV/LATE for EFW on macro outcomes using randomization inference.

\textbf{Assumptions.}
\begin{itemize}[leftmargin=1.5em]
\item Within $\abs{m}\le r$, assignment is as-if random (stronger than continuity, but local).
\item Exclusion and relevance as before, but now in finite-sample experimental framing.
\end{itemize}

\textbf{Failure modes (skeptical referee).}
\begin{itemize}[leftmargin=1.5em]
\item Choosing $r$ can be subjective; small $r$ yields small samples; large $r$ violates as-if random.
\item Cross-country heterogeneity makes a single $r$ less defensible; must allow system-specific windows.
\end{itemize}

\textbf{Robustness/sensitivity/bounds.}
\begin{enumerate}[leftmargin=1.8em]
\item Window selection based on covariate balance with transparent rules; show robustness across $r$ grid.
\item Use exact Fisher randomization tests for reduced form and first stage; then invert to get IV confidence sets.
\end{enumerate}

\vspace{0.8em}

%%%%%%%%%%%%%%%%%%%%%%%%%%%%%%%%%%%%%%%%%%%%%%%%%%%%%%%%%%%%%%%%%%%%%%%%%%%%%%%
\subsection*{5.4 Self-critique as skeptical referee (Part 2)}

\textbf{Referee critique (hard-nosed).}
\begin{itemize}[leftmargin=1.5em]
\item Strategy 1 (ratio) is transparent and avoids some parametric pitfalls, but it still relies
heavily on exclusion.
\item Strategy 2 (2SLS) is flexible and can incorporate multiple EFW components, but invites
specification and bandwidth disputes.
\item Strategy 3 (local randomization) is appealing for finite samples but hinges on window choice
and may lack power for tail outcomes.
\end{itemize}

\textbf{Most serious threat across all strategies:} exclusion/direct effects of $D$ on $Y$ not through EFW.

Therefore, the proposal must:
\begin{itemize}[leftmargin=1.5em]
\item treat exclusion as partially falsifiable via auxiliary tests,
\item quantify robustness to violations via bounds,
\item and avoid claiming global ``institutions $\rightarrow$ macro'' effects beyond local compliers.
\end{itemize}

\vspace{0.8em}

%%%%%%%%%%%%%%%%%%%%%%%%%%%%%%%%%%%%%%%%%%%%%%%%%%%%%%%%%%%%%%%%%%%%%%%%%%%%%%%
\subsection*{5.5 Default choice and final implementable specification}

\textbf{Default strategy: hybrid RD-IV anchored by robust RD estimation (Strategy 1+2).}\
We choose Strategy 1 (local Wald ratio) as the headline estimand for transparency,
implemented using robust bias-corrected RD estimates for numerator and denominator. ([rdpackages.github.io][10])\
We use Strategy 2 (2SLS within bandwidth) as a supportive implementation that enables:
\begin{itemize}[leftmargin=1.5em]
\item multiple endogenous EFW components,
\item interactions for state dependence,
\item and nonlinear/asymmetric models requiring multiple instruments.
\end{itemize}

We treat Strategy 3 (local randomization) as a pre-registered robustness appendix for key outcomes.

\vspace{0.4em}

\textbf{Final implementable RD-IV specification (baseline).}\
For horizon $h$, estimate reduced form and first stage via robust RD:
\begin{align}
\hat{\tau}_{Y}(h)
&=
\widehat{\RD}\big(\Delta Y_{ce}(h), m_{ce}, D_{ce}\big),
\\
\hat{\tau}_{E}(h)
&=
\widehat{\RD}\big(\Delta \EFW_{ce}(h), m_{ce}, D_{ce}\big),
\end{align}
and report:
\begin{equation}
\hat{\beta}^{\IV}(h)=\hat{\tau}_{Y}(h)/\hat{\tau}_{E}(h),
\end{equation}
with confidence intervals computed by delta method and by weak-IV robust inversion (appendix).

\vspace{0.4em}

\textbf{Exclusion sensitivity as a first-class deliverable.}\
We define a ``direct winner effect'' parameter $\delta(h)$:
\[
\Delta Y_{ce}(h) = \beta(h)\Delta \EFW_{ce}(h) + \delta(h)D_{ce} + \tilde{u}_{ce}(h).
\]
If $\delta(h)\ne 0$, the Wald ratio is biased.
We will:
\begin{itemize}[leftmargin=1.5em]
\item calibrate plausible bounds on $\delta(h)$ from placebo outcomes and short-run policy-irrelevant outcomes,
\item compute identified sets for $\beta(h)$ under $\delta(h)\in[\ubarDelta(h),\barDelta(h)]$,
\item and present these bounds alongside point estimates.
\end{itemize}

\vspace{1.0em}

%%%%%%%%%%%%%%%%%%%%%%%%%%%%%%%%%%%%%%%%%%%%%%%%%%%%%%%%%%%%%%%%%%%%%%%%%%%%%%%
\section*{6) REQUIRED MODULE A --- Asymmetry of EFW effects}
\textit{(EFW$\uparrow$ vs EFW$\downarrow$; $\ge 3$ designs; self-critique; pick default; implementable spec.)}

\vspace{0.4em}

%%%%%%%%%%%%%%%%%%%%%%%%%%%%%%%%%%%%%%%%%%%%%%%%%%%%%%%%%%%%%%%%%%%%%%%%%%%%%%%
\subsection*{6.1 Three asymmetry designs (identification-safe)}

\textbf{Goal:} identify whether macro responses differ for EFW increases (reforms) versus decreases
(reversals), without selecting on post-treatment sign.

\vspace{0.4em}

\textbf{Design A1: Two-endogenous-regressor asymmetric IV with predetermined interaction instruments.}\
\begin{itemize}[leftmargin=1.5em]
\item \textbf{Estimand:} $(\beta^+(h),\beta^-(h))$ in \eqref{eq:asymstruct}.
\item \textbf{Identification idea:} instrument $\Delta \EFW^+_{ce}(h)$ and $\Delta \EFW^-_{ce}(h)$ with
$\bZ_{ce}=(D_{ce}, D_{ce}\times S_{ce})$, where $S_{ce}$ is predetermined and chosen to shift the
probability/magnitude of reversals vs reforms (e.g., high baseline EFW creates ``room to reverse'').
\item \textbf{Exact hypotheses:}
\[
H_0: \beta^+(h)=\beta^-(h)
\qquad
\text{vs}
\qquad
H_1: \beta^+(h)\ne \beta^-(h).
\]
\item \textbf{Comparison:} slope of $Y$ on positive EFW changes vs slope on negative EFW changes.
\item \textbf{Assumptions:}
\begin{itemize}[leftmargin=1.5em]
\item RD continuity and instrument exogeneity (local).
\item Relevance: instruments must independently predict $\Delta \EFW^+$ and $\Delta \EFW^-$.
\item Exclusion: instruments affect $Y$ only via EFW components.
\end{itemize}
\item \textbf{Pitfalls:}
\begin{itemize}[leftmargin=1.5em]
\item Weak instruments for one component (especially reversals) can make asymmetry hard to identify.
\item Exclusion may fail differentially by sign (e.g., left winners might use fiscal channels not in EFW).
\end{itemize}
\item \textbf{Sellable outputs:} table of $\hat{\beta}^+(h),\hat{\beta}^-(h)$ and test of equality;
figure showing IV-predicted $\Delta \EFW^+$ and $\Delta \EFW^-$ by instrument variation.
\end{itemize}

\vspace{0.5em}

\textbf{Design A2: Pre-determined `reform-prone'' vs `reversal-prone'' regimes with interaction-IV.}\
\begin{itemize}[leftmargin=1.5em]
\item \textbf{Estimand:} differential IV effect by predetermined regime $R_{ce}\in{0,1}$:
\[
\beta(h;R=0),\ \beta(h;R=1),
\quad
\text{and test }
H_0:\beta(h;R=0)=\beta(h;R=1).
\]
\item \textbf{Regime definition (pre-election only):}
$R_{ce}=1$ if baseline EFW is high (above pre-specified threshold) and the contest features a
credible less-market challenger (large negative ideology gap), making reversals plausible;
$R_{ce}=0$ otherwise (reforms more plausible).
\item \textbf{Method:} estimate IV with interactions:
\[
\Delta Y = \beta_0 \Delta \EFW + \beta_1 (\Delta \EFW\times R) + u,
\]
instrumenting $(\Delta \EFW, \Delta \EFW\times R)$ with $(D, D\times R)$.
\item \textbf{Assumptions:} same as IV + regime exogeneity (predetermined).
\item \textbf{Pitfalls:} regime classification may proxy for other channels (e.g., high EFW countries differ).
\item \textbf{Sellable outputs:} interaction plot of $\hat{\beta}(h)$ across regimes; pre-trend checks by regime.
\end{itemize}

\vspace{0.5em}

\textbf{Design A3: Asymmetric local Wald via separate ``reversal-encouragement'' instrument.}\
\begin{itemize}[leftmargin=1.5em]
\item \textbf{Estimand:} LATE for reversal channel vs reform channel using two instruments:
\[
Z^{\reform}_{ce}=D_{ce}\times \1{\EFW_{c,t_{ce}-1}\le \kappa},
\qquad
Z^{\reversal}_{ce}=(1-D_{ce})\times \1{\EFW_{c,t_{ce}-1}\ge \kappa},
\]
where $\kappa$ is a pre-specified baseline EFW threshold.
\item \textbf{Idea:} in low-EFW environments, a more-market win is more likely to generate reform;
in high-EFW environments, a less-market win is more likely to generate reversal.
\item \textbf{Implementation:} two-sample IV (or interacted 2SLS) to estimate slopes in reform vs reversal
encouragement environments and compare.
\item \textbf{Assumptions/pitfalls:}
\begin{itemize}[leftmargin=1.5em]
\item Requires that these instruments have nontrivial first stages in corresponding subsamples.
\item Still vulnerable to exclusion differences by regime.
\end{itemize}
\item \textbf{Sellable outputs:} first-stage plots showing where reversals/reforms occur; asymmetry in macro impacts.
\end{itemize}

\vspace{0.8em}

%%%%%%%%%%%%%%%%%%%%%%%%%%%%%%%%%%%%%%%%%%%%%%%%%%%%%%%%%%%%%%%%%%%%%%%%%%%%%%%
\subsection*{6.2 Self-critique as skeptical referee (asymmetry)}

\textbf{Referee critique.}
\begin{itemize}[leftmargin=1.5em]
\item A1 is the cleanest direct test of $\beta^+\neq\beta^-$, but it relies on having sufficiently
strong and distinct instruments for positive and negative EFW movements.
\item A2 is feasible and naturally integrates with state dependence, but it risks re-labeling
`state dependence'' as `asymmetry'' unless the regime truly predicts sign changes.
\item A3 is conceptually appealing (separate encouragements) but may be weak if reversals are rare
or if high-EFW countries have different macro dynamics unrelated to EFW changes.
\end{itemize}

\textbf{Key empirical feasibility concern:} are election-induced EFW decreases sufficiently common?
This is testable in Part 1 by estimating first-stage distributions and sign frequencies.

\vspace{0.8em}

%%%%%%%%%%%%%%%%%%%%%%%%%%%%%%%%%%%%%%%%%%%%%%%%%%%%%%%%%%%%%%%%%%%%%%%%%%%%%%%
\subsection*{6.3 Default asymmetry design and evidence package}

\textbf{Default: Design A1 (instrumented asymmetric piecewise effect), with A2 as robustness.}\
We choose A1 because it directly targets the asymmetry estimand in \eqref{eq:asymstruct}
and provides a clean equality test.
We use A2 as a robustness framework to interpret asymmetry through predetermined regimes.

\vspace{0.5em}

\textbf{Final implementable specification (core equations).}\
Define endogenous regressors:
\[
E^+_{ce}(h)=\Delta \EFW^+_{ce}(h),
\qquad
E^-_{ce}(h)=\Delta \EFW^-_{ce}(h).
\]

Second stage:
\begin{equation}
\Delta Y_{ce}(h)
=
\alpha(h)
+
\beta^+(h)E^+_{ce}(h)
+
\beta^-(h)E^-_{ce}(h)
+
g(m_{ce})
+
u_{ce}(h),
\label{eq:A1second}
\end{equation}
with instruments:
\begin{equation}
\bZ_{ce}
=
\big(D_{ce}, D_{ce}\times S_{ce}\big),
\label{eq:A1instr}
\end{equation}
where $S_{ce}$ is predetermined and chosen from a small pre-specified set that is most likely to
shift reform vs reversal capacity:
\begin{itemize}[leftmargin=1.5em]
\item baseline EFW level $\EFW_{c,t_{ce}-1}$ (room to reform vs room to reverse),
\item baseline inflation history (high-inflation states constrain policy responses),
\item baseline institutional constraints (V-Dem constraints on executive). ([v-dem.net][18])
\end{itemize}

We estimate \eqref{eq:A1second} within RD bandwidth and report weak-IV robust inference.

\vspace{0.4em}

\textbf{Exact hypothesis tests.}
\begin{itemize}[leftmargin=1.5em]
\item Equality: $H_0:\beta^+(h)=\beta^-(h)$ via Wald test with weak-IV robust p-value.
\item Directional: $H_0:\beta^+(h)\le \beta^-(h)$ where theory suggests reversals hurt more (or vice versa).
\item Joint across horizons: pre-specified family-wise error control (e.g., Holm) for $h\in{1,2,4}$.
\end{itemize}

\vspace{0.4em}

\textbf{Headline tables/figures (3--5).}
\begin{enumerate}[leftmargin=1.8em]
\item Table A: first-stage sign distribution of $\Delta \EFW(h)$ by winner near cutoff.
\item Table B: $\hat{\beta}^+(h),\hat{\beta}^-(h)$ for growth and investment; equality tests.
\item Figure A: binned scatter of $\Delta \EFW(h)$ vs $D$ by baseline EFW state (shows instrument relevance).
\item Figure B: asymmetric IV impulse response (horizon plot) for $\beta^+(h)$ and $\beta^-(h)$.
\item Table C: sensitivity/bounds for exclusion (range of direct effect $\delta$) on asymmetry conclusions.
\end{enumerate}

\vspace{0.4em}

\textbf{Required robustness/sensitivity/bounds.}
\begin{itemize}[leftmargin=1.5em]
\item Weak-IV robust confidence sets for $(\beta^+,\beta^-)$.
\item Alternative instrument sets (swap $S_{ce}$ choices) with pre-specified order.
\item Exclusion sensitivity: allow direct effects $\delta(h)$ and show how large $\delta$ must be to remove
asymmetry.
\item Re-estimate using EFW areas excluding mechanically overlapping components for inflation/tail outcomes.
\end{itemize}

\vspace{1.0em}

%%%%%%%%%%%%%%%%%%%%%%%%%%%%%%%%%%%%%%%%%%%%%%%%%%%%%%%%%%%%%%%%%%%%%%%%%%%%%%%
\section*{7) REQUIRED MODULE B --- State dependence of EFW effects}
\textit{(EFW$\rightarrow$macro state dependence, not merely winner effects;
$\ge 3$ designs; self-critique; choose default; implementable.)}

\vspace{0.4em}

%%%%%%%%%%%%%%%%%%%%%%%%%%%%%%%%%%%%%%%%%%%%%%%%%%%%%%%%%%%%%%%%%%%%%%%%%%%%%%%
\subsection*{7.1 Three state-dependence designs}

\textbf{Design B1: Interaction-IV within RD bandwidth (continuous state).}\
\begin{itemize}[leftmargin=1.5em]
\item \textbf{Estimand:} $\beta_0(h)$ and $\beta_1(h)$ in:
\[
\Delta Y = \beta_0 \Delta \EFW + \beta_1 (\Delta \EFW \times S) + u,
\]
interpretable as marginal effect varying with $S$.
\item \textbf{Instruments:} $(D, D\times S)$ for $(\Delta\EFW, \Delta\EFW\times S)$.
\item \textbf{Assumptions:} RD continuity; exclusion; relevance of both instruments.
\item \textbf{Pitfalls:} if $S$ is highly correlated with omitted channels of $D$ on $Y$, exclusion is shakier.
\item \textbf{Sellable outputs:} plot of implied $\beta(h;S)$ across $S$ with confidence bands.
\end{itemize}

\vspace{0.5em}

\textbf{Design B2: Stratified local-IV (CATE-IV) by predetermined bins of state.}\
\begin{itemize}[leftmargin=1.5em]
\item \textbf{Estimand:} $\beta^{\IV}(h;S\in B_j)$ for bins $B_j$ of $S$ (e.g., terciles of baseline EFW).
\item \textbf{Method:} run RD-IV separately within each bin; report heterogeneity tests.
\item \textbf{Assumptions:} within-bin RD continuity; enough close elections within each bin.
\item \textbf{Pitfalls:} loss of power; multiple testing; risk of ``bin mining'' if bins not pre-specified.
\item \textbf{Sellable outputs:} heterogeneity table and binned coefficient plot.
\end{itemize}

\vspace{0.5em}

\textbf{Design B3: Local randomization + IV within matched windows.}\
\begin{itemize}[leftmargin=1.5em]
\item \textbf{Estimand:} finite-sample LATE in a small window for each state regime.
\item \textbf{Method:} choose window $r$ by balance; then within window, estimate IV and perform randomization inference.
\item \textbf{Assumptions:} as-if random within window; exclusion.
\item \textbf{Pitfalls:} window choice and small samples in each regime; may be feasible only for few outcomes.
\item \textbf{Sellable outputs:} exact p-values for heterogeneity of EFW effect by state.
\end{itemize}

\vspace{0.8em}

%%%%%%%%%%%%%%%%%%%%%%%%%%%%%%%%%%%%%%%%%%%%%%%%%%%%%%%%%%%%%%%%%%%%%%%%%%%%%%%
\subsection*{7.2 Self-critique as skeptical referee (state dependence)}

\textbf{Referee critique.}
\begin{itemize}[leftmargin=1.5em]
\item B1 is efficient and interpretable, but hinges on strong exclusion with interactions.
\item B2 is transparent but can be underpowered and vulnerable to multiple testing.
\item B3 is compelling when sample is small, but may not scale across countries/horizons.
\end{itemize}

\vspace{0.4em}

\textbf{Multiplicty risk is central.}
State dependence is an invitation to specification search.
We must pre-register a small set of states, justify them theoretically, and enforce a multiple-testing plan.

\vspace{0.8em}

%%%%%%%%%%%%%%%%%%%%%%%%%%%%%%%%%%%%%%%%%%%%%%%%%%%%%%%%%%%%%%%%%%%%%%%%%%%%%%%
\subsection*{7.3 Default state-dependence design}

\textbf{Default: B1 (interaction-IV) with pre-specified states; B2 as robustness visualization.}\

\vspace{0.4em}

\textbf{Primary state variables (pre-election only) and rationale.}\
We pre-specify a short list ($\le 5$) of state variables:

\begin{enumerate}[leftmargin=1.8em]
\item \textbf{Baseline EFW level:} $\EFW_{c,t_{ce}-1}$.
Rationale: ``room to reform'' and institutional complementarities; also predicts reversal feasibility.
\item \textbf{Inflation regime history:} mean and volatility of inflation in $t_{ce}-5$ to $t_{ce}-1$.
Rationale: policy constraints and credibility issues; may shape returns to institutional change.
\item \textbf{Institutional constraints / democracy:} V-Dem indices (e.g., constraints on executive).
Rationale: ability to implement reforms; credibility of property-rights changes. ([v-dem.net][18])
\item \textbf{Openness:} trade openness (WDI) and baseline GDP per capita.
Rationale: gains from trade/market reforms plausibly larger in open and/or poorer economies.
\item \textbf{State capacity proxy:} pre-election governance effectiveness or tax capacity (where available).
Rationale: enforcement and implementation capacity for legal/regulatory reforms.
\end{enumerate}

\vspace{0.4em}

\textbf{Exact tests and plots.}
\begin{itemize}[leftmargin=1.5em]
\item For each state $S$, estimate interaction-IV and test $H_0:\beta_1(h)=0$.
\item Plot implied $\beta(h;S)$ over the interquartile range of $S$.
\item Present ``state dependence map'': two-dimensional heatmap for (baseline EFW, inflation history)
using coarse bins (pre-specified) to avoid overfitting.
\end{itemize}

\vspace{0.4em}

\textbf{Multiple testing discipline.}
\begin{itemize}[leftmargin=1.5em]
\item Pre-specify the list of states, horizons, and outcomes for heterogeneity.
\item Use Holm adjustment within each outcome family (growth, inflation, investment, tails).
\item Emphasize effect sizes and confidence intervals, not only p-values.
\end{itemize}

\vspace{1.0em}

%%%%%%%%%%%%%%%%%%%%%%%%%%%%%%%%%%%%%%%%%%%%%%%%%%%%%%%%%%%%%%%%%%%%%%%%%%%%%%%
\section*{8) REQUIRED MODULE C --- Nonlinearities of EFW effects}
\textit{(Nonlinearities in EFW$\rightarrow$macro; $\ge 3$ designs; self-critique; default; pre-spec plan.)}

\vspace{0.4em}

%%%%%%%%%%%%%%%%%%%%%%%%%%%%%%%%%%%%%%%%%%%%%%%%%%%%%%%%%%%%%%%%%%%%%%%%%%%%%%%
\subsection*{8.1 Three nonlinearity designs}

\textbf{Design C1: Instrumented threshold (piecewise slope) in baseline EFW.}\
\begin{itemize}[leftmargin=1.5em]
\item \textbf{Estimand:} different slopes below vs above threshold $\kappa$:
\[
\Delta Y
=
\beta_L \Delta \EFW\cdot \1\{\EFW_0\le \kappa\}
+
\beta_H \Delta \EFW\cdot \1\{\EFW_0> \kappa\}
+
u.
\]
\item \textbf{Instruments:} $D\cdot \1\{\EFW_0\le \kappa\}$ and $D\cdot \1\{\EFW_0> \kappa\}$
for the two endogenous interacted regressors.
\item \textbf{Assumptions:} RD continuity; exclusion; relevance in both regimes.
\item \textbf{Pitfalls:}
\begin{itemize}[leftmargin=1.5em]
\item If $\kappa$ is chosen post hoc, p-hacking risk is extreme.
\item If one regime has weak first stage, results unstable.
\end{itemize}
\item \textbf{Deliverables:} piecewise IV slope plot and test $\beta_L=\beta_H$.
\end{itemize}

\vspace{0.5em}

\textbf{Design C2: Instrumented magnitude bins of $\Delta \EFW$ (nonlinear dose response).}\
\begin{itemize}[leftmargin=1.5em]
\item \textbf{Estimand:} effect differs for `small'' vs `large'' EFW changes.
\item \textbf{Implementation:} define endogenous dose indicators:
\[
B^{(1)}=\1{0<\Delta \EFW \le q_{50}},
\quad
B^{(2)}=\1{\Delta \EFW > q_{50}},
\]
and interact with $\Delta \EFW$ to build piecewise linear response.
\item \textbf{Identification:} use predetermined-interaction instruments that shift magnitude
(e.g., $D\times$ baseline EFW, $D\times$ institutional constraints) to generate variation in dose.
\item \textbf{Pitfalls:} high dimensionality and weak instruments; needs strict pre-specification.
\item \textbf{Deliverables:} estimated IV marginal effects by dose; monotonicity diagnostics.
\end{itemize}

\vspace{0.5em}

\textbf{Design C3: Tail-focused nonlinearities (EFW effect on downside risk).}\
\begin{itemize}[leftmargin=1.5em]
\item \textbf{Estimand:} effect of EFW on tail outcomes, e.g., probability of large drawdowns or inflation spikes:
\[
\Prob(\mathrm{MDD}\ge d^\star \mid \Delta\EFW),
\quad
\Prob(\mathrm{InflSpike}=1 \mid \Delta\EFW).
\]
\item \textbf{Implementation:} IV linear probability model for tail indicators and IV quantile-type summaries
for continuous downside measures (worst growth, drawdown magnitude).
\item \textbf{Assumptions:} exclusion; interpretability of LPM IV; robustness to functional form.
\item \textbf{Pitfalls:} tail events are rare; power concerns; sensitive to data quality.
\item \textbf{Deliverables:} IV estimates for tail indicators and drawdown severity; horizon plots.
\end{itemize}

\vspace{0.8em}

%%%%%%%%%%%%%%%%%%%%%%%%%%%%%%%%%%%%%%%%%%%%%%%%%%%%%%%%%%%%%%%%%%%%%%%%%%%%%%%
\subsection*{8.2 Self-critique as skeptical referee (nonlinearities)}

\textbf{Referee critique.}
\begin{itemize}[leftmargin=1.5em]
\item C1 is interpretable and feasible, but threshold choice must be disciplined.
\item C2 is ambitious and risks weak instruments; may be too complex for credible identification.
\item C3 is policy-relevant and distinct (tails), but power is the binding constraint and must be emphasized.
\end{itemize}

Given credibility constraints, a conservative default should focus on C1 and C3,
with C2 relegated to exploratory appendix if instruments are strong enough.

\vspace{0.8em}

%%%%%%%%%%%%%%%%%%%%%%%%%%%%%%%%%%%%%%%%%%%%%%%%%%%%%%%%%%%%%%%%%%%%%%%%%%%%%%%
\subsection*{8.3 Default nonlinearity design and pre-specification plan}

\textbf{Default: C1 (baseline EFW threshold) + C3 (tail outcomes), pre-specified.}\

\vspace{0.4em}

\textbf{Threshold choice plan (avoid fishing).}
\begin{itemize}[leftmargin=1.5em]
\item Pre-specify $\kappa$ as the median baseline EFW in the estimation sample (or fixed percentile),
computed from pre-election EFW only (no outcome information).
\item Alternatively, use two thresholds at terciles to show monotone patterns (but adjust multiplicity).
\item Report robustness to small shifts in $\kappa$ (e.g., $\pm 0.25$ EFW points) as a sensitivity plot.
\end{itemize}

\vspace{0.4em}

\textbf{Exact hypotheses and tests.}
\begin{itemize}[leftmargin=1.5em]
\item $H_0:\beta_L(h)=\beta_H(h)$ (no threshold nonlinearity).
\item $H_0:\beta_{\tail}(h)=0$ for each tail outcome family.
\item Joint across horizons with Holm correction.
\end{itemize}

\vspace{0.4em}

\textbf{Core plots.}
\begin{itemize}[leftmargin=1.5em]
\item Piecewise IV slope plot for $\beta_L(h)$ and $\beta_H(h)$.
\item Horizon plot for tail effects: probability of inflation spike, probability of large drawdown.
\item ``Nonlinearity diagnostic'': first-stage strength by baseline EFW regime.
\end{itemize}

\vspace{1.0em}

%%%%%%%%%%%%%%%%%%%%%%%%%%%%%%%%%%%%%%%%%%%%%%%%%%%%%%%%%%%%%%%%%%%%%%%%%%%%%%%
\section*{9) Novelty}
\textit{(8--12 citations total; 1 sentence each; minimal essential foundations.)}

\vspace{0.4em}

\textbf{Reference set (12 items; minimal and sharp).}\
We will cite these as the backbone of the paper:

\begin{enumerate}[leftmargin=1.8em]
\item Lee (2008): foundational close-election RD logic and diagnostic implication that baseline covariates
should be continuous at the cutoff. ([ScienceDirect][11])
\item Calonico, Cattaneo, and Titiunik (2014): robust bias-corrected RD inference and bandwidth logic
used for primary RD estimation. ([rdpackages.github.io][10])
\item Cattaneo, Frandsen, and Titiunik (2015): local randomization inference in RD as a finite-sample robustness
check. ([rdpackages.github.io][13])
\item Imbens and Angrist (1994): LATE framework clarifying what IV identifies (complier-local effects).
\item Conley, Hansen, and Rossi (2012): ``plausibly exogenous'' IV sensitivity analysis to bound causal effects
when exclusion is imperfect.
\item Gwartney, Lawson, and Murphy (2025): EFW dataset definition, coverage, components, and official citation
guidance. ([Fraser Institute][1])
\item CLEA project: global constituency-level election returns enabling construction of close-election margins
in a broad cross-country setting. ([electiondataarchive.org][2])
\item ParlGov project: clean parliamentary election and cabinet data for EU/OECD, enabling a high-precision
subset and validation sample. ([parlgov.org][3])
\item World Bank DPI: broad cross-country political institutions and ideology variables to classify
market orientation where finer sources are unavailable. ([Inter-American Development Bank][4])
\item V-Dem / V-Party: expert-coded party positions and institutional states for orientation and state dependence,
with documented dataset access and versioning. ([v-dem.net][5])
\item World Bank WDI: primary macro outcomes via open API and CC BY 4.0 licensing for replication. ([datahelpdesk.worldbank.org][7])
\item Penn World Table: macro levels and productivity measures, with open licensing and recent version
covering years through 2023. ([University of Groningen][8])
\end{enumerate}

\vspace{1.0em}

%%%%%%%%%%%%%%%%%%%%%%%%%%%%%%%%%%%%%%%%%%%%%%%%%%%%%%%%%%%%%%%%%%%%%%%%%%%%%%%
\section*{10) Threats / Failure Modes}
\textit{(Brutally honest; diagnostics + mitigation + fallback + claims if it fails.)}

\vspace{0.4em}

\textbf{Threat 1: Exclusion/direct effects (Part 2).}\
\begin{itemize}[leftmargin=1.5em]
\item \textbf{Problem:} winner affects macro outcomes through channels not captured by EFW.
\item \textbf{Diagnostics:}
\begin{itemize}[leftmargin=1.5em]
\item test RD effects on outcomes plausibly orthogonal to EFW (short-run administrative outcomes),
\item check component-mediation plausibility: if macro effects are large but EFW changes are tiny,
exclusion is doubtful,
\item compare results when using EFW subsets excluding components that may embed the outcome.
\end{itemize}
\item \textbf{Mitigation:} sensitivity/bounds for direct effects; interpret as `effect of the EFW-induced
policy bundle'' rather than pure EFW.
\item \textbf{Fallback:} if bounds are too wide, downgrade Part 2 to `mediation-consistent'' evidence:
report reduced-form and first-stage, and present IV as exploratory with transparent caveats.
\item \textbf{Claim if it fails:} ``close-election market orientation affects macro outcomes and EFW
moves, but mapping to a causal EFW effect is not identified without stronger assumptions.''
\end{itemize}

\vspace{0.5em}

\textbf{Threat 2: Weak instrument / weak first stage.}\
\begin{itemize}[leftmargin=1.5em]
\item \textbf{Problem:} $\tau_{\Delta \EFW(h)}$ may be small for some horizons or contexts.
\item \textbf{Diagnostics:} report first-stage estimates, effective sample sizes, and weak-IV robust tests.
\item \textbf{Mitigation:} focus on horizons where first stage is strongest; emphasize reduced form otherwise.
\item \textbf{Fallback:} component-wise first stages (some areas may move more than overall EFW).
\item \textbf{Claim if it fails:} Part 2 becomes infeasible for causal EFW; Part 1 remains publishable.
\end{itemize}

\vspace{0.5em}

\textbf{Threat 3: Measurement error in EFW and ideology score.}\
\begin{itemize}[leftmargin=1.5em]
\item \textbf{Problem:} EFW is a composite with evolving sources; ideology labels may be noisy.
\item \textbf{Diagnostics:} replicate Part 1 first stage across alternative ideology sources (V-Party vs DPI vs MARPOR).
\item \textbf{Mitigation:} pre-specify ideology pipeline; drop ambiguous contests; use continuous ideology gap.
\item \textbf{Fallback:} treat $D$ as ``winner ideology'' instrument and interpret estimates as effect of
market-oriented governance rather than precise EFW shifts.
\item \textbf{Claim if it fails:} reduced-form winner effects remain; EFW channel interpretation weakens.
\end{itemize}

\vspace{0.5em}

\textbf{Threat 4: Sign-specific instrument validity (asymmetry).}\
\begin{itemize}[leftmargin=1.5em]
\item \textbf{Problem:} identifying $\beta^+$ and $\beta^-$ requires instruments that separately move
positive and negative EFW components; may be weak or invalid.
\item \textbf{Diagnostics:} show first-stage decomposition into $\Delta \EFW^+$ and $\Delta \EFW^-$ by instruments;
report weak-IV diagnostics.
\item \textbf{Mitigation:} limit asymmetry claims to cases with strong first-stage decomposition; otherwise
present as suggestive.
\item \textbf{Fallback:} asymmetry in reduced form (winner effects) remains valid even if EFW asymmetry fails.
\item \textbf{Claim if it fails:} ``no credible evidence on asymmetric EFW slopes''; retain state dependence
in reduced form and/or in EFW changes.
\end{itemize}

\vspace{0.5em}

\textbf{Threat 5: Cross-country comparability and external validity.}\
\begin{itemize}[leftmargin=1.5em]
\item \textbf{Problem:} RD estimates are local and may vary by institutions; pooling may mask heterogeneity.
\item \textbf{Diagnostics:} subgroup analyses by region, regime type, baseline democracy, electoral system.
\item \textbf{Mitigation:} interpret as local effects; avoid global structural claims.
\item \textbf{Fallback:} present a `high-precision'' subsample (e.g., ParlGov democracies) as main,
and global extension as supplementary.
\item \textbf{Claim if it fails:} `effects are context-specific; the paper maps heterogeneity rather than
a single global parameter.''
\end{itemize}

\vspace{1.0em}

%%%%%%%%%%%%%%%%%%%%%%%%%%%%%%%%%%%%%%%%%%%%%%%%%%%%%%%%%%%%%%%%%%%%%%%%%%%%%%%
\section*{11) Paper Output Plan}
\textit{(10--14 core tables/figures + appendix checklist.)}

\vspace{0.4em}

\textbf{Core tables/figures (proposed 14).}

\begin{enumerate}[leftmargin=1.8em]
\item Figure 1: RD plot of $\Delta \EFW(h)$ at cutoff for $h=1$ (first stage, overall EFW).
\item Figure 2: RD plots for EFW areas (5 small panels shown sequentially; report separately in paper/appendix).
\item Table 1: First-stage estimates $\hat{\tau}_{\Delta \EFW(h)}$ across horizons $h\in{0,1,2,4}$.
\item Table 2: Reduced-form macro effects $\hat{\tau}_{\Delta Y(h)}$ for growth, inflation, investment.
\item Figure 3: Horizon plot of reduced-form macro effects (term-aligned).
\item Table 3: Reduced-form tail outcomes: worst growth, inflation spike, max drawdown, crisis start.
\item Figure 4: RD-IV (local Wald) estimates $\hat{\beta}^{\IV}(h)$ for growth and investment (horizon plot).
\item Table 4: RD-IV estimates and weak-IV robust confidence sets for key horizons.
\item Table 5: Exclusion sensitivity bounds (direct effect $\delta$ grid) for $\beta^{\IV}(h)$.
\item Figure 5: Sensitivity plot: identified set width vs assumed direct effect magnitude.
\item Table 6 (Asymmetry): $\hat{\beta}^+(h),\hat{\beta}^-(h)$ and equality tests (Module A).
\item Figure 6 (Asymmetry): asymmetric IV impulse response: reforms vs reversals.
\item Table 7 (State dependence): interaction-IV coefficients for baseline EFW, inflation history, institutions.
\item Figure 7 (Nonlinearity): piecewise IV slopes by baseline EFW threshold + tail effects plot.
\end{enumerate}

\vspace{0.6em}

\textbf{Appendix diagnostics checklist (pre-specified).}
\begin{itemize}[leftmargin=1.5em]
\item Density test at cutoff; report heaping/rounding sensitivity.
\item Covariate balance at cutoff for predetermined covariates (levels and pre-trends).
\item Donut RD sensitivity ($\eta$ grid).
\item Placebo cutoffs ($\pm 1\%,\pm 2\%$) and placebo outcomes (pre-election).
\item Bandwidth sensitivity and polynomial order sensitivity.
\item Election-type stratification (presidential vs parliamentary).
\item Ideology-scoring robustness (V-Party vs DPI vs MARPOR).
\item EFW-component overlap robustness (exclude Sound Money for inflation outcomes).
\item Weak-IV robust inference (Anderson--Rubin inversion) for all IV claims.
\item Exclusion sensitivity/bounds for all IV claims (report identified sets).
\end{itemize}

\vspace{1.5em}

%%%%%%%%%%%%%%%%%%%%%%%%%%%%%%%%%%%%%%%%%%%%%%%%%%%%%%%%%%%%%%%%%%%%%%%%%%%%%%%
\appendix

%%%%%%%%%%%%%%%%%%%%%%%%%%%%%%%%%%%%%%%%%%%%%%%%%%%%%%%%%%%%%%%%%%%%%%%%%%%%%%%
\section*{Appendix A: Concrete implementation notes (non-ETL; feasibility-focused)}

\textbf{A.1 Expected sample size and feasibility logic (conservative).}\
Even with global election coverage, the usable RD sample will be a strict subset:
\begin{itemize}[leftmargin=1.5em]
\item must have reliable close-election margins,
\item must have non-ambiguous executive ``winner'' mapping,
\item must have credible ideology orientation,
\item must have EFW and macro data for required horizons,
\item must pass basic balance and manipulation diagnostics.
\end{itemize}

Therefore, feasibility strategy is \emph{tiered}:
\begin{itemize}[leftmargin=1.5em]
\item Tier 1: ParlGov democracies (high quality, clean mapping, strong internal validity).
\item Tier 2: Expanded sample using CLEA + auxiliary mapping for executive formation.
\item Tier 3: Global breadth via DPI ideology, with transparent coarseness and robustness emphasis.
\end{itemize}

\textbf{A.2 Why we expect a first stage from close-election winners to EFW.}\
EFW is designed to reflect economic policy institutions (size of government, legal system, regulation,
trade openness, and sound money). ([Fraser Institute][1])\
Market-oriented winners plausibly change:
\begin{itemize}[leftmargin=1.5em]
\item trade policy (tariffs, quotas, capital controls),
\item regulatory policy (business/labor regulation),
\item privatization and size-of-government dimensions,
\item legal/property rights enforcement (longer horizon).
\end{itemize}
The first stage is empirical and must be shown; we do not assume it.

\textbf{A.3 A note on using inflation outcomes.}\
Because inflation enters EFW ``Sound Money'', inflation-as-outcome must avoid mechanical overlap. ([Fraser Institute][1])\
We implement:
\begin{itemize}[leftmargin=1.5em]
\item inflation outcomes paired with EFW excluding Sound Money, and/or
\item EFW areas excluding inflation-related subcomponents, and
\item separate analysis of Sound Money as an endogenous policy/outcome bundle (interpretation cautious).
\end{itemize}

\vspace{1.0em}

%%%%%%%%%%%%%%%%%%%%%%%%%%%%%%%%%%%%%%%%%%%%%%%%%%%%%%%%%%%%%%%%%%%%%%%%%%%%%%%
\section*{Appendix B: Pre-analysis plan style pre-specification (discipline against fishing)}

\textbf{B.1 Primary outcomes (locked).}\
\begin{itemize}[leftmargin=1.5em]
\item Growth: real GDP per capita growth (WDI and PWT versions).
\item Investment: gross fixed capital formation \% GDP (WDI) and investment share (PWT where applicable).
\item Inflation: CPI inflation (WDI), but only interpreted with EFW excluding Sound Money.
\item Tail: max drawdown in log GDP per capita, worst growth, inflation spike ($\pi^\star \in {20,40}$),
and crisis start (Laeven--Valencia 1970--2011 where available). ([SSRN][9])
\end{itemize}

\textbf{B.2 Primary horizons (locked).}\
\begin{itemize}[leftmargin=1.5em]
\item $h\in{1,2,4}$ as primary (short, medium, within-term).
\item $h=0$ and $h=5$ as secondary (immediate and longer).
\item Term-aligned averages as additional primary for comparability.
\end{itemize}

\textbf{B.3 Primary EFW measures (locked).}\
\begin{itemize}[leftmargin=1.5em]
\item Overall EFW.
\item 5 areas as separate endogenous measures.
\item ``EFW excluding Sound Money'' for inflation outcomes (constructed as average of remaining areas).
\end{itemize}

\textbf{B.4 Primary RD specification (locked).}\
\begin{itemize}[leftmargin=1.5em]
\item Local linear, triangular kernel, robust bias-corrected inference, MSE-optimal bandwidth. ([rdpackages.github.io][10])
\item Country-clustered standard errors (wild bootstrap as robustness).
\end{itemize}

\textbf{B.5 Primary Part-2 IV specification (locked).}\
\begin{itemize}[leftmargin=1.5em]
\item Wald ratio for headline $\beta^{\IV}(h)$.
\item 2SLS within bandwidth for models with interactions/piecewise terms.
\item Weak-IV robust inference and exclusion sensitivity bounds always reported for IV.
\end{itemize}

\textbf{B.6 State variables (locked list, $\le 5$).}\
\begin{itemize}[leftmargin=1.5em]
\item Baseline EFW.
\item Inflation history (mean/volatility).
\item V-Dem constraints on executive (or democracy index). ([v-dem.net][18])
\item Trade openness.
\item Baseline log GDP per capita.
\end{itemize}

\textbf{B.7 Nonlinearity thresholds (locked).}\
\begin{itemize}[leftmargin=1.5em]
\item Baseline EFW median threshold as default.
\item Secondary: terciles with Holm correction.
\end{itemize}

\textbf{B.8 Asymmetry instruments (locked order).}\
We pre-specify $S_{ce}$ candidate list in \eqref{eq:A1instr} in this order:
\begin{enumerate}[leftmargin=1.8em]
\item baseline EFW,
\item baseline inflation history,
\item V-Dem constraints,
\item trade openness,
\item baseline log GDP per capita.
\end{enumerate}
We stop after the first two that generate sufficiently strong first-stage decomposition
(weak-IV diagnostics guide reporting but do not change specification search beyond this list).

\vspace{1.0em}

%%%%%%%%%%%%%%%%%%%%%%%%%%%%%%%%%%%%%%%%%%%%%%%%%%%%%%%%%%%%%%%%%%%%%%%%%%%%%%%
\section*{Appendix C: Detailed falsification and robustness menu (implementable checklist)}

\textbf{C.1 RD validity checks (mandatory).}\
\begin{itemize}[leftmargin=1.5em]
\item Density test at $m=0$; report separately by election type and by region.
\item Balance tests on predetermined covariates; report joint F-test.
\item Pre-trend placebo: RD on $\EFW_{c,t_{ce}-1}-\EFW_{c,t_{ce}-2}$ and on pre-election growth.
\item Donut RD excluding $\eta \in {0.1,0.25,0.5}$ percentage points (or comparable).
\item Placebo cutoffs: $m= \pm 1\%, \pm 2\%$.
\end{itemize}

\textbf{C.2 Outcome definition robustness.}\
\begin{itemize}[leftmargin=1.5em]
\item Growth from WDI vs PWT; show both.
\item Investment from WDI gross fixed capital formation vs PWT investment share.
\item Inflation outcomes interpreted only with EFW excluding Sound Money.
\item Tail thresholds sensitivity: $\pi^\star$ alternatives and drawdown threshold $d^\star$ alternatives.
\end{itemize}

\textbf{C.3 Orientation robustness.}\
\begin{itemize}[leftmargin=1.5em]
\item Use V-Party vs DPI vs MARPOR labeling where overlapping.
\item Use continuous ideology gap as alternative to binary more/less-market labeling.
\item Exclude ambiguous ideology-gap elections ($|I_A-I_B|<\delta_I$) and report sensitivity.
\end{itemize}

\textbf{C.4 Inference robustness.}\
\begin{itemize}[leftmargin=1.5em]
\item Country-clustered vs two-way clustered (country-year) where applicable.
\item Wild cluster bootstrap for small number of clusters.
\item Local randomization inference window results for key outcomes.
\end{itemize}

\textbf{C.5 Exclusion sensitivity (mandatory for Part 2).}\
\begin{itemize}[leftmargin=1.5em]
\item Report identified sets for $\beta(h)$ under bounded direct effect $\delta(h)$.
\item Calibrate $\delta(h)$ bounds from placebo outcomes and short-run outcomes not plausibly affected by EFW.
\item Component-wise IV: compare results for Legal System and Regulation vs Size of Government and Trade.
\end{itemize}

\vspace{1.0em}

%%%%%%%%%%%%%%%%%%%%%%%%%%%%%%%%%%%%%%%%%%%%%%%%%%%%%%%%%%%%%%%%%%%%%%%%%%%%%%%
\section*{Appendix D: Data-access quick reference (provider, constraints, and replication posture)}

\textbf{EFW.}\
Dataset downloadable for free; report provides formal citation and notes variable/source evolution. ([Fraser Institute][1])

\textbf{Elections.}\
CLEA provides constituency-level returns for many elections and is described as publicly available at no cost,
though some access may require registration; versioning exists. ([electiondataarchive.org][2])

\textbf{ParlGov.}\
Open data infrastructure for EU/OECD democracies 1900--2023; downloadable releases. ([parlgov.org][3])

\textbf{Political institutions / ideology.}\
DPI2020 covers 1975--2020 and includes ideology measures; used as fallback. ([Inter-American Development Bank][4])

\textbf{Party positions.}\
V-Dem provides datasets for download; V-Party is used for party economic positions. ([v-dem.net][5])

\textbf{Manifesto data.}\
Manifesto Project access may require API key/registration; manifestoR documents API retrieval. ([CRAN][6])

\textbf{Macro outcomes.}\
WDI via open API with CC BY 4.0 licensing; PWT downloads with CC BY 4.0 (10.0) and newer version 11.0
through 2023. ([datahelpdesk.worldbank.org][7])

\textbf{Crises.}\
Laeven--Valencia update provides crisis chronology 1970--2011; used where overlapping with sample. ([SSRN][9])

\vspace{1.0em}

%%%%%%%%%%%%%%%%%%%%%%%%%%%%%%%%%%%%%%%%%%%%%%%%%%%%%%%%%%%%%%%%%%%%%%%%%%%%%%%
\section*{Appendix E: Minimal bibliography skeleton (paper will include full BibTeX)}

\textbf{Note:} This proposal uses a minimal reference set (Section 9) and will keep the paper’s
literature review short and identification-focused.

\vspace{0.6em}

\textbf{(End of proposal.)}

\end{document}
